\documentclass[10pt, letterpaper]{article}

% Inhaltsverzeichnis für Pakettypen (nur für Übersicht im Header, wird nicht im Dokument angezeigt)
% 1. Seitenlayout und Ränder
% 2. Sprache und Zeichensatz
% 3. Mathematik und Theorem-Umgebungen
% 4. Eigene Makros
% 5. Diagramme und Grafiken
% 6. Tabellen und Aufzählungen
% 7. Inhaltsverzeichnis
% 8. Abschnittsüberschriften
% 9. Abstrakt-Umgebung
% 10. Todos/Notizen
% 11. Rahmen/Box-Umgebungen
% 12. Python-Integration
% 13. Literaturverwaltung
% 14. Hyperlinks
% 15. Absatzeinstellungen
% 16. Umgebungen
% 17  Graphik
% 18  Extra
% 00. Titel und Autor

\usepackage{slashed}

% --- 1. Seitenlayout und Ränder ---
\usepackage[margin=3cm]{geometry}

% --- 2. Sprache und Zeichensatz ---
\usepackage[english]{babel}
\usepackage[T1]{fontenc}
\usepackage[utf8]{inputenc}

% --- 3. Mathematik und Theorem-Umgebungen ---
\usepackage{amsmath, amssymb, amsthm}
\usepackage{mathrsfs}
\DeclareMathOperator{\WF}{WF}

% --- 4. Eigene Makros ---
\usepackage{xcolor}
\newcommand{\SKP}{\langle\cdot,\cdot\rangle}
\newcommand{\id}{\operatorname{id}}
\newcommand{\R}{\mathbb{R}}
\newcommand{\N}{\mathbb{N}}
\newcommand{\Q}{\mathbb{Q}}
\newcommand{\Z}{\mathbb{Z}}
\newcommand{\C}{\mathbb{C}}
\newcommand{\QU}{\mathbb{H}}
\newcommand{\Cinf}{C^{\infty}}
\newcommand{\Mod}{\mathrm{Mod}}
\newcommand{\Alg}{\mathrm{Alg}}
\newcommand{\Aut}{\mathrm{Aut}}
\newcommand{\End}{\mathrm{End}}
\newcommand{\Hom}{\mathrm{Hom}}
\newcommand{\Cl}{\mathrm{Cl}}
\newcommand{\entwurf}[1]{\textcolor{red}{#1}}
\newcommand{\GL}{\mathrm{GL}}
\newcommand{\Mat}{\mathrm{Mat}}
\newcommand{\SL}{\mathrm{SL}}
\newcommand{\SO}{\mathrm{SO}}
\newcommand{\SU}{\mathrm{SU}}
\newcommand{\Sp}{\mathrm{Sp}}
\newcommand{\Spin}{\mathrm{Spin}}
\newcommand{\Ogroup}{\mathrm{O}}
\newcommand{\U}{\mathrm{U}}

% --- 5. Diagramme und Grafiken ---
\usepackage{graphicx}
\graphicspath{{../../Library/Images/}}
\usepackage[export]{adjustbox}
\usepackage{tikz}
\usetikzlibrary{decorations.pathreplacing, arrows.meta, positioning}
\usepackage{tikz-cd}

% --- 6. Tabellen und Aufzählungen ---
\usepackage{enumitem}
\setlist[itemize]{left=0.5cm}

\newenvironment{romanenum}[1][]
  {%
    \ifx&#1&
    \else
      \textbf{#1}\quad
    \fi
    \begin{enumerate}[label=\roman*)]
  }
  {%
    \end{enumerate}%
  }

% --- 7. Inhaltsverzeichnis ---
\usepackage{tocloft}
\renewcommand{\cftsecfont}{\footnotesize}
\renewcommand{\cftsubsecfont}{\footnotesize}
\renewcommand{\cftsubsubsecfont}{\footnotesize}
\renewcommand{\cftsecpagefont}{\footnotesize}
\renewcommand{\cftsubsecpagefont}{\footnotesize}
\renewcommand{\cftsubsubsecpagefont}{\footnotesize}
\usepackage{etoc}

% --- 8. Abschnittsüberschriften ---
\usepackage{titlesec}
\titleformat{\section}{\normalfont\large\bfseries}{\thesection}{1em}{}
\titleformat{\subsection}{\normalfont\normalsize\bfseries}{\thesubsection}{0.5em}{}
\titleformat{\subsubsection}{\normalfont\normalsize\bfseries}{\thesubsubsection}{0.5em}{}
\setcounter{secnumdepth}{4}

% --- 9. Abstrakt-Umgebung ---
\usepackage{changepage}
\renewenvironment{abstract}
  {
    \begin{adjustwidth}{1.5cm}{1.5cm}
    \small
    \textsc{Abstract. –}%
  }
  {
    \end{adjustwidth}
  }

% --- 10. Todos/Notizen ---
\usepackage{todonotes}
% Hellblau definieren
\definecolor{hellblau}{rgb}{0.3,0.6,1}
\newenvironment{note}{\par\begingroup\color{hellblau}\itshape}{\par\endgroup}

% --- 11. Rahmen/Box-Umgebungen ---
\usepackage{mdframed}
\usepackage{tcolorbox}
\colorlet{shadecolor}{gray!25}

\newenvironment{customTheorem}
  {\vspace{10pt}%
   \begin{mdframed}[
     backgroundcolor=gray!20,
     linewidth=0pt,
     innertopmargin=10pt,
     innerbottommargin=10pt,
     skipabove=\dimexpr\topsep+\ht\strutbox\relax,
     skipbelow=\topsep,
   ]}
  {\end{mdframed}
   \vspace{10pt}%
  }

% --- 12. Python-Integration ---
% (Deaktiviert in dieser Version, aktiviere bei Bedarf)
% \usepackage{pythontex}
% \usepackage[makestderr]{pythontex}

% --- 13. Literaturverwaltung ---
\usepackage{csquotes}
\usepackage[backend=biber, style=alphabetic, citestyle=alphabetic]{biblatex}
\addbibresource{bibliography.bib}

% --- 14. Hyperlinks ---
\usepackage{hyperref}
\hypersetup{
  colorlinks   = true,
  urlcolor     = blue,
  linkcolor    = blue,
  citecolor    = blue,
  frenchlinks  = true
}

% --- 15. Absatzeinstellungen ---
\usepackage[parfill]{parskip}
\sloppy

% --- 16. Umgebungen ---
\usepackage{thmtools}

\newcommand{\CustomHeading}[3]{%
  \par\medskip\noindent%
  \textbf{#1 #2} \textnormal{(#3)}.\enskip%
}

\newenvironment{DEF}[2]{\begin{unitbox}\CustomHeading{Definition}{#1}{#2}}{\end{unitbox}}
\newenvironment{PROP}[2]{\begin{unitbox}\CustomHeading{Proposition}{#1}{#2}}{\end{unitbox}}
\newenvironment{THEO}[2]{\begin{unitbox}\CustomHeading{Theorem}{#1}{#2}}{\end{unitbox}}
\newenvironment{LEM}[2]{\begin{unitbox}\CustomHeading{Lemma}{#1}{#2}}{\end{unitbox}}
\newenvironment{KORO}[2]{\begin{unitbox}\CustomHeading{Corollar}{#1}{#2}}{\end{unitbox}}
\newenvironment{REM}[2]{\begin{unitbox}\CustomHeading{Remark}{#1}{#2}}{\end{unitbox}}
\newenvironment{EXA}[2]{\begin{unitbox}\CustomHeading{Example}{#1}{#2}}{\end{unitbox}}
\newenvironment{STUD}[2]{\begin{unitbox}\CustomHeading{Study}{#1}{#2}}{\end{unitbox}}
\newenvironment{CONC}[2]{\begin{unitbox}\CustomHeading{Concept}{#1}{#2}}{\end{unitbox}}
\newenvironment{OTH}[2]{\begin{unitbox}\CustomHeading{Other}{#1}{#2}}{\end{unitbox}}
\newenvironment{EXE}[2]{\begin{unitbox}\CustomHeading{Exercise}{#1}{#2}}{\end{unitbox}}
\newenvironment{MOT}[2]{\begin{unitbox}\CustomHeading{Motivation}{#1}{#2}}{\end{unitbox}}
\newenvironment{PROOF}[2]{\begin{unitbox}\CustomHeading{Proof}{#1}{#2}}{\end{unitbox}}

% --- Unit Umgebung für Source-Inhalte ---
\usepackage{mdframed}
\newmdenv[
  linewidth=1pt,
  topline=false,
  bottomline=false,
  rightline=false,
  leftmargin=0cm,
  rightmargin=0cm,
  skipabove=10pt,
  skipbelow=10pt,
  innertopmargin=0.5\baselineskip,
  innerbottommargin=0.5\baselineskip,
  backgroundcolor=gray!10,
  linecolor=gray
]{unitbox}

\newenvironment{unit}[1]
  {\begin{unitbox}\textbf{Unit #1}\par\smallskip}
  {\end{unitbox}}

% --- 17. ... ---

% --- 18. Extras ---
\usepackage{stmaryrd}
\usepackage{bbold}  % falls du \mathbb{1} nutzen willst

% --- 00. Titel und Autor ---
\title{Mein Titel}
\author{Tim Jaschik}
\date{\today}

\begin{document}

\maketitle
\rule{\textwidth}{0.5pt}
\begin{abstract}
Kurze Beschreibung …
\end{abstract}
\rule{\textwidth}{0.5pt}
\vspace{0.5cm}

\tableofcontents

\pagebreak



\begin{abstract}
Wir geben einen top-down Einstieg in Algebren über einem kommutativen Ring \(R\).
Ausgehend von der allgemeinen Definition beschreiben wir die freie (Tensor-)Algebra und die daraus
durch Relationen entstehenden klassischen Beispiele (symmetrische und äußere Algebren),
diskutieren gradierte Strukturen und Division-Algebren und platzieren das Ganze in einen
kategorientheoretischen Rahmen (\emph{Monoidenobjekte in \(\Mod_R\)} und Adjunktionen).
\end{abstract}

\section{Top-Down: Was ist eine \(R\)-Algebra?}

\begin{DEF}[Algebra über einem Ring]
Sei \(R\) ein \emph{kommutativer} Ring mit \(1\).
Eine \emph{\(R\)-Algebra} ist ein \(R\)-Modul \(A\) zusammen mit einer \(R\)-bilinearen Multiplikation
\[
\mu:A\times A\to A,\quad (a,b)\mapsto ab,
\]
so dass \((A,+,\mu,1_A)\) ein assoziativer Ring mit Eins ist und die Skalarmultiplikation aus \(R\) mit \(\mu\) verträglich ist:
\[
r\cdot(ab)=(ra)b=a(rb)\qquad(\forall\,r\in R,\;a,b\in A).
\]
Ein Homomorphismus \(f:A\to B\) von \(R\)-Algebren ist ein \(R\)-linearer Ringhomomorphismus.
\end{DEF}

\begin{REM}[Schnittstelle Ring/Modul; kategoriale Sicht]
Die Kategorie \(\Alg_R\) der \(R\)-Algebren ist die Kategorie der \emph{Monoidenobjekte} in \(\Mod_R\):
Eine \(R\)-Algebra ist genau ein Objekt \(A\in\Mod_R\) mit Multiplikation \(\mu:A\otimes_R A\to A\) und Einheit \(\eta:R\to A\),
die die Monoidenaxiome erfüllen.
Der \emph{Vergissfunktor} \(U:\Alg_R\to\Mod_R\) besitzt einen linken Adjunkten \(T:\Mod_R\to\Alg_R\)
(die \emph{freie Algebra}, siehe §\ref{sec:tensor-algebra}).
\end{REM}

\begin{EXA}[Grundbeispiele]
\begin{enumerate}[label=(\alph*), leftmargin=*]
\item \textbf{Polynomringe:} \(R[x_1,\dots,x_n]\) ist eine \(R\)-Algebra via koeffizientenweise \(R\)-Linearität.
\item \textbf{Matrizenalgebren:} \(M_n(R)\) ist eine \(R\)-Algebra mit Matrixmultiplikation.
\item \textbf{Endomorphismen:} \(\mathrm{End}_R(M)\) ist eine \(R\)-Algebra für jeden \(R\)-Modul \(M\).
\item \textbf{Gruppenalgebra:} Für endliche Gruppe \(G\) ist \(R[G]\) die freie \(R\)-Modulsumme über \(G\) mit Faltung.
\end{enumerate}
\end{EXA}

\section{Universelle Bilinearität: Tensorprodukt und Linearisierung}
\begin{DEF}[Tensorprodukt]
Für \(R\)-Moduln \(M,N\) ist das \emph{Tensorprodukt} \(M\otimes_R N\) ein \(R\)-Modul
zusammen mit einer bilinearen Abbildung \(\tau:M\times N\to M\otimes_R N\) mit folgender \emph{universeller Eigenschaft}:
Für jeden \(R\)-Modul \(P\) induziert jede bilineare Abbildung \(f:M\times N\to P\) eindeutig einen \(R\)-linearen
Homomorphismus \(\tilde f:M\otimes_R N\to P\) mit \(\tilde f(m\otimes n)=f(m,n)\).
\end{DEF}

\begin{REM}[Warum das Tensorprodukt zentral ist]
Bilinearität der Multiplikation \(\mu:A\times A\to A\) einer \(R\)-Algebra entspricht eindeutig einem
\(R\)-linearen Morphismus \(\bar\mu:A\otimes_R A\to A\).
Damit sind Algebren genau die Objekte \(A\) in \(\Mod_R\) ausgestattet mit \(\bar\mu\) und \(\eta:R\to A\),
die die Monoidenaxiome erfüllen.
\end{REM}

\section{Freie Algebren und universelle Konstruktionen}
\label{sec:tensor-algebra}
\begin{DEF}[Tensoralgebra / freie \(R\)-Algebra]
Für einen \(R\)-Modul \(V\) heißt
\[
T_R(V)\;:=\;\bigoplus_{n=0}^{\infty} V^{\otimes n}
\quad(\text{mit }V^{\otimes 0}:=R)
\]
\emph{Tensoralgebra} von \(V\), mit der durch Konkatenation induzierten Multiplikation
\(V^{\otimes i}\otimes V^{\otimes j}\to V^{\otimes(i+j)}\).
Sie ist die \emph{freie \(R\)-Algebra} über \(V\): für jede \(R\)-Algebra \(A\) und jeden \(R\)-linearen
\(f:V\to A\) gibt es genau einen Algebra\-homomorphismus \(\Phi:T_R(V)\to A\) mit \(\Phi|_V=f\).
\end{DEF}

\begin{REM}[Quotienten durch Relationen]
Viele klassische Algebren entstehen als Quotienten \(T_R(V)/\langle\mathcal{R}\rangle\) einer freien Algebra
durch ein \emph{homogenes} Relatorideal \(\langle\mathcal{R}\rangle\). Beispiele folgen in §\ref{sec:symm-ext}.
\end{REM}

\section{Symmetrische und äußere Algebra}
\label{sec:symm-ext}
\begin{DEF}[Symmetrische Algebra]
Sei \(V\) ein \(R\)-Modul. Die \emph{symmetrische Algebra} ist
\[
S_R(V):=T_R(V)\big/\langle v\otimes w - w\otimes v \mid v,w\in V\rangle.
\]
Sie ist die freie \emph{kommutative} \(R\)-Algebra, die \(V\) enthält: Für jede kommutative \(R\)-Algebra \(A\)
und \(R\)-lineares \(f:V\to A\) gibt es genau einen kommutativen Algebrahomomorphismus \(S_R(V)\to A\) mit
der gewünschten Restriktion.
\end{DEF}

\begin{DEF}[]
Die \emph{äußere Algebra} eines \(R\)-Moduls \(V\) ist
\[
\Lambda_R(V):=T_R(V)\big/\langle v\otimes v \mid v\in V\rangle.
\]
Sie ist die freie \emph{alternierende} Algebra über \(V\), d.\,h.\ das Produkt \(\wedge\) erfüllt \(v\wedge v=0\).
\end{DEF}

\begin{REM}[Graduierung]
Sowohl \(S_R(V)\) als auch \(\Lambda_R(V)\) sind \emph{natürlich gradierte} Algebren:
\[
S_R(V)=\bigoplus_{n\ge 0} S^n_R(V),\qquad
\Lambda_R(V)=\bigoplus_{n\ge 0} \Lambda^n_R(V),
\]
mit \(S^n_R(V)\cdot S^m_R(V)\subseteq S^{n+m}_R(V)\) bzw.
\(\Lambda^n_R(V)\wedge \Lambda^m_R(V)\subseteq \Lambda^{n+m}_R(V)\).
\end{REM}

\begin{EXA}[]
Über einem Körper \(k\) ist \(S_k(V)\) kanonisch isomorph zum Polynomring auf \(V^\ast\) (endliche Dimension),
während \(\Lambda_k(V)\) die Algebra der alternierenden Tensoren bildet (Grundlage der Differentialformen).
\end{EXA}

\section{Graduierte Algebren}
\begin{DEF}[Gradierte Algebra]
Eine \emph{(insbesondere \(\N\)-)gradierte \(R\)-Algebra} ist eine \(R\)-Algebra \(A=\bigoplus_{n\in\Z}A_n\)
mit \(A_i\cdot A_j\subseteq A_{i+j}\).
Ist \(A\) kommutativ, spricht man oft von \emph{super-}gradierter Struktur via Parität \(A_{\bar 0}\oplus A_{\bar 1}\).
\end{DEF}

\begin{REM}[Relevanz]
Gradierungen sind strukturerhaltend: Homomorphismen respektieren die Grade;
in der Topologie tragen Kohomologieringe \(H^\ast(X;R)\) eine natürliche Gradierung,
in der Linearen Algebra sind \(S^n(V),\Lambda^n(V)\) die homogenen Komponenten.
\end{REM}

\section{Division-Algebren}
\begin{DEF}[Division-\(R\)-Algebra]
Eine (nicht notwendig kommutative) \(R\)-Algebra \(D\) heißt \emph{Division-Algebra},
wenn jedes \(0\neq d\in D\) ein (links- und rechts-)inverses Element besitzt.
Ist \(R\) ein Körper \(k\), spricht man von \emph{Division-\(k\)-Algebra}.
\end{DEF}

\begin{REM}
Klassische Beispiele über \(k=\mathbb{R}\) sind \(\mathbb{R},\mathbb{C},\mathbb{H}\) (Hurwitz).
Division-Struktur ist \emph{keine} Eigenschaft freier Konstruktionen; sie ist restriktiv
und eng mit Darstellungstheorie und quadratischen Formen verknüpft.
\end{REM}

\section{Kategorientheoretische Kommentare (Leitplanken)}
\begin{itemize}[leftmargin=*]
\item \textbf{Monoidenobjekte:} \(\Alg_R\simeq \mathrm{Mon}(\Mod_R)\). Die Multiplikation ist ein Morphismus
\(\bar\mu:A\otimes_R A\to A\); die Einheit ist \(\eta:R\to A\).
\item \textbf{Adjunktion \((T\dashv U)\):} Der freie-Vergiss-Adjunktionszusammenhang
\[
\Hom_{\Alg_R}(T_R(V),A)\;\cong\;\Hom_{\Mod_R}(V,\,U(A))
\]
charakterisiert \(T_R(V)\) universell.
\item \textbf{Quotienten durch Relationen:} Symmetrische und äußere Algebra sind
linke adjungierte zu \emph{symmetrischen} bzw.\ \emph{alternierenden} Funktoren
(Erzwingen von \(vw=wv\) bzw.\ \(v\wedge v=0\)).
\item \textbf{Gradierungen:} Gradierte Algebren sind Monoidenobjekte in der abelschen Kategorie
\(\Mod_R^{\Z}\) mit monoidaler Struktur \((A\otimes B)_n=\bigoplus_{i+j=n}A_i\otimes B_j\).
\end{itemize}


\section{Clifford-Algebren}
\label{sec:clifford}

\begin{DEF}[Clifford-Algebra]
Sei \(R\) ein kommutativer Ring und \(V\) ein \(R\)-Modul, versehen mit einer
symmetrischen \(R\)-bilinearen Form
\[
Q:V\times V\to R.
\]
Die \emph{Clifford-Algebra} \(\Cl(V,Q)\) ist definiert als der Quotient der Tensoralgebra
\[
\Cl(V,Q):=T_R(V)\big/\left\langle v\otimes v - Q(v)\,1 \mid v\in V \right\rangle.
\]
\end{DEF}

\begin{PROP}[Universelle Eigenschaft]
Für jede \(R\)-Algebra \(A\) und jeden \(R\)-linearen Homomorphismus
\(f:V\to A\) mit der Eigenschaft \(f(v)^2=Q(v)\,1_A\)
existiert genau ein \(R\)-Algebrahomomorphismus
\(\Phi:\Cl(V,Q)\to A\) mit \(\Phi(v)=f(v)\).
\end{PROP}

\begin{REM}[Beziehung zu Tensor- und äußerer Algebra]
\begin{itemize}[leftmargin=*]
\item Für \(Q=0\) gilt \(\Cl(V,0)\cong \Lambda_R(V)\), da die Relation \(v^2=0\) der antisymmetrischen
Relation \(v\wedge v=0\) entspricht.
\item Für nicht-triviales \(Q\) deformiert \(\Cl(V,Q)\) die äußere Algebra:
\[
v\wedge w + w\wedge v = 2\,Q(v,w)\,1.
\]
\end{itemize}
\end{REM}

\begin{EXA}[Klassische Clifford-Algebren]
Über \(R=\mathbb{R}\) und einer positiv-definiten Form \(Q(x)=\|x\|^2\) entstehen:
\[
\Cl_0(\mathbb{R})=\mathbb{R},\quad
\Cl_1(\mathbb{R})\simeq\mathbb{C},\quad
\Cl_2(\mathbb{R})\simeq\mathbb{H}.
\]
Diese Beispiele führen zu den bekannten Bott-periodischen Reihen
\(\Cl_{p+8,q}\cong \Cl_{p,q}\otimes\Mat_{16}(\mathbb{R})\).
\end{EXA}

\begin{REM}[Graduierung]
\(\Cl(V,Q)\) ist eine \(\mathbb{Z}_2\)-\emph{graduierte Algebra}:
\[
\Cl(V,Q)=\Cl^0(V,Q)\oplus \Cl^1(V,Q),
\]
wobei \(\Cl^0\) der Unterraum der Produkte gerader Länge und \(\Cl^1\) der ungerader Länge ist.
Diese Parität bildet die Grundlage für \emph{Superalgebren}.
\end{REM}

\section{Superalgebren und Superkommutativität}

\begin{DEF}[Superalgebra]
Eine \emph{Superalgebra} über \(R\) ist eine \(\mathbb{Z}_2\)-graduierte Algebra
\[
A=A_{\bar{0}}\oplus A_{\bar{1}}
\]
mit Multiplikation
\(A_i\cdot A_j\subseteq A_{i+j \bmod 2}\).
\end{DEF}

\begin{DEF}[Superkommutativität]
Eine Superalgebra heißt \emph{superkommutativ}, wenn
\[
ab = (-1)^{|a||b|}\,ba
\qquad \text{für alle homogenen }a,b\in A,
\]
wobei \(|a|\in\{0,1\}\) den Grad bezeichnet.
\end{DEF}

\begin{EXA}[Superkommutative Beispiele]
\begin{itemize}[leftmargin=*]
\item Die äußere Algebra \(\Lambda_R(V)\) ist superkommutativ,
da \(v\wedge w = - w\wedge v\).
\item Die Kohomologiealgebra \(H^\ast(X;R)\) ist superkommutativ:
\[
\alpha\smile\beta = (-1)^{pq}\,\beta\smile\alpha
\quad (\alpha\in H^p,\,\beta\in H^q).
\]
\end{itemize}
\end{EXA}

\begin{REM}[Kategorischer Standpunkt]
Superalgebren sind Monoidenobjekte in der monoidalen Kategorie
\((\Mod_R^{\mathbb{Z}_2},\otimes_R)\),
deren Tauschisomorphismus das Vorzeichen \((-1)^{|a||b|}\) trägt.
Dies motiviert den Begriff der \emph{superkommutativen} Monoide.
\end{REM}


\pagebreak

\section{Der algebraische Zugang zur reellen Clifford-Algebra nach Baker (Kap.~5)}
\label{sec:baker-clifford}

\begin{MOT}{Algebraische Sichtweise}{mot-baker}
Bakers Darstellung der Clifford-Algebren (Kapitel~5 in \cite{Baker2002})
folgt einem strikt algebraischen Prinzip:
eine Clifford-Algebra ist ein \emph{freies assoziatives Objekt} mit
Relationen, die eine gegebene bilineare Form kodieren.
Die universelle Eigenschaft erlaubt dann, lineare Abbildungen
$f:V\to A$ mit $f(v)^2=-Q(v)$ auf eindeutige Algebrahomomorphismen
$F:\Cl(V,Q)\to A$ zu erweitern.
\end{MOT}






\label{subsec:baker-to-quotient}

\begin{MOT}{Ziel}{mot-baker-to-quotient}
Wir zeigen, dass Bakers Präsentation der reellen Clifford-Algebren über Erzeugern
$e_1,\dots,e_n$ mit Relationen
\[
e_i^2=-1,\qquad e_ie_j=-e_je_i\ (i\ne j)
\]
\emph{äquivalent} ist zur koordinatenfreien Definition als Quotient der freien
assoziativen Algebra $T(V)$ durch das von der Form $Q$ induzierte Relatorideal.
\end{MOT}

\begin{DEF}{Freie Algebra und kanonische Einbettung}{def-free-embed}
Sei $V$ ein endlichdimensionaler $\R$-Vektorraum und $T(V)=\bigoplus_{k\ge0}V^{\otimes k}$ die Tensoralgebra
mit Konkatenationsmultiplikation. Die kanonische Inklusion bezeichne
\[
j:V\hookrightarrow T(V),\qquad v\mapsto v\in V^{\otimes 1}\subset T(V).
\]
\end{DEF}

\begin{DEF}{Quadratische Form, Polarisierung, Relationen}{def-Q-B-rel}
Sei $Q:V\to\R$ eine quadratische Form mit \emph{polarer Form}
\[
B(v,w)\;=\;\tfrac12\big(Q(v+w)-Q(v)-Q(w)\big)
\quad(\text{symmetrisch, bilinear}).
\]
Wir definieren das (zweiseitige) \emph{Clifford-Ideal} $I_Q\unlhd T(V)$ als das kleinste Ideal, das die Elemente
\[
j(v)\otimes j(v)+Q(v)\cdot 1 \quad (v\in V)
\]
enthält (Äquivalenzformen s.\ Remark~\ref{rem:equiv-rel}).
\end{DEF}

\begin{DEF}{Clifford-Algebra als Tensorquotient}{def-cl-quot}
Die \emph{Clifford-Algebra} zu $(V,Q)$ ist der Quotient
\[
\Cl(V,Q)\;:=\;T(V)\big/I_Q.
\]
Wir schreiben das Produkt in $\Cl(V,Q)$ als $\cdot$ oder (kontextabhängig) punktlos.
Die Kompositionsabbildung $V\xrightarrow{j}T(V)\twoheadrightarrow \Cl(V,Q)$
bezeichnen wir mit $\iota:V\to\Cl(V,Q)$.
\end{DEF}

\begin{REM}{Äquivalente Präsentationen der Relationen}
\label{rem:equiv-rel}
Die Definition über $j(v)^2+Q(v)1=0$ ist äquivalent zu der \emph{bilinearen} Relation
\[
\iota(v)\iota(w)+\iota(w)\iota(v)\;=\;2\,B(v,w)\,1\qquad(\forall v,w\in V),
\]
denn Polarsummen $Q(v+w)$ expandieren genau zu $B(v,w)$.
Umgekehrt folgt die quadratische Relation aus der bilinearen durch Einsetzen $v=w$.
\end{REM}

\begin{PROP}{Bakers Präsentation als Spezialfall der Quotient-Definition}{prop:baker-as-special}
Sei $(V,Q)=(\R^n,\|\cdot\|^2)$ und wähle eine $Q$-orthonormale Basis $(e_1,\dots,e_n)$.
Dann hat $\Cl(V,Q)$ die Präsentation
\[
\Cl_n\;=\;\langle\,e_1,\dots,e_n\mid e_i^2=-1,\ e_ie_j=-e_je_i\ (i\ne j)\,\rangle,
\]
also genau Bakers Generator-Relationen.
\end{PROP}

\begin{PROOF}{Skizze}{proof:baker-as-special}
Für eine $Q$-ONB gilt $B(e_i,e_j)=\delta_{ij}$ und $Q(e_i)=1$.
Damit liefert Remark~\ref{rem:equiv-rel}:
\[
e_ie_j+e_je_i=2\delta_{ij}\,1.
\]
Für $i=j$ ist das $2e_i^2=2\cdot 1\cdot(-1)$, also $e_i^2=-1$ (unser Vorzeichenkonvent:
$\iota(v)^2=-Q(v)1$). Für $i\ne j$ folgt $e_ie_j=-e_je_i$.
Damit sind genau Bakers Relationen erfüllt; umgekehrt erzeugen diese Relationen
das in Definition~\ref{def-cl-quot} eingeführte Ideal $I_Q$.
\end{PROOF}

\begin{PROP}{Universelle Eigenschaft: Quotient $\Longleftrightarrow$ Baker}{prop:univ-eq}
Die folgenden Aussagen sind äquivalent:
\begin{enumerate}[label=(\alph*), leftmargin=*]
\item \textbf{(Quotient-Universalität)} Für jede $\R$-Algebra $A$ und lineare Abbildung
$f:V\to A$ mit $f(v)^2=-Q(v)\,1_A$ existiert genau ein Algebrahomomorphismus
$\Phi:\Cl(V,Q)\to A$ mit $\Phi\circ\iota=f$.
\item \textbf{(Baker-Universalität)} Sei $A$ eine $\R$-Algebra mit Elementen
$F_1,\dots,F_n$ (Bakers \emph{Generatoren}) mit Relationen
$F_i^2=-1$ und $F_iF_j=-F_jF_i$ ($i\ne j$). Dann gibt es genau einen
Algebrahomomorphismus $\Phi:\Cl_n\to A$ mit $\Phi(e_i)=F_i$.
\end{enumerate}
\end{PROP}

\begin{PROOF}{Skizze}{proof:univ-eq}
(a)$\Rightarrow$(b): Wähle $V=\R^n$ mit ONB $(e_i)$ und $f(e_i)=F_i$.
Dann $f(e_i)^2=-1$ und $f(e_if_j+f_je_i)=0$ für $i\ne j$,
also existiert $\Phi$ eindeutig.
(b)$\Rightarrow$(a): Gegeben $f:V\to A$ mit $f(v)^2=-Q(v)$.
Wähle ONB $(e_i)$ und setze $F_i=f(e_i)$; dann erfüllt $(F_i)$ Bakers Relationen
und induziert $\Phi:\Cl_n\to A$. Unabhängigkeit von der ONB folgt aus der
Koordinatenfreiheit der Relation $\iota(v)^2=-Q(v)1$.
\end{PROOF}

\begin{THEO}{Basis und Filterung (PBW für Clifford)}{thm:pbw-cl}
Mit einer $Q$-ONB $(e_1,\dots,e_n)$ besitzt $\Cl(V,Q)$ die
$\R$-Vektorraum-Basis
\[
\{\,e_{i_1}\cdots e_{i_r}\mid 0\le r\le n,\ 1\le i_1<\dots<i_r\le n\,\},
\]
also $\dim_\R\Cl(V,Q)=2^n$.
\end{THEO}

\begin{PROOF}{Idee}{proof:pbw-cl}
Benutze die aufsteigende Filterung $F^r\Cl$ durch Wortlängen und den
assoziierten Graduiertenring $\mathrm{gr}(\Cl)$. Die Relation
$v^2+Q(v)1=0$ ist vom Grad $\le2$ und verschwindet im führenden Grad,
so dass $\mathrm{gr}(\Cl)\cong\Lambda(V)$. Damit heben sich
die üblichen monomischen Basisargumente (wie bei der äußeren Algebra)
auf $\Cl$ zurück; dies liefert die behauptete Basis.
\end{PROOF}

\begin{REM}{Signaturen und allgemeine Konventionen}{rem:signatures}
Für eine Signatur $(p,q)$ mit einer ONB $(e_1,\dots,e_{p+q})$ und
\[
Q(e_i)=\begin{cases}
+1,& 1\le i\le p,\\
-1,& p<i\le p+q,
\end{cases}
\]
lautet die Relation koordiniert
\[
e_i^2=-Q(e_i)\,1=\begin{cases}
-1,& 1\le i\le p,\\
+1,& p<i\le p+q,
\end{cases}
\qquad
e_ie_j=-e_je_i\ (i\ne j),
\]
also Bakers Präsentation mit gemischten Quadratsignen. Koordinatenfrei bleibt stets
$\iota(v)^2=-Q(v)1$ und $\iota(v)\iota(w)+\iota(w)\iota(v)=2B(v,w)1$.
\end{REM}

\begin{CONC}{Matrix-Modelle via Universalität}{conc:matrix-models}
Um ein \emph{konkretes} Modell $\Cl(V,Q)\to A$ (z.\,B. $A=M_{2^k}(\R),M_{2^k}(\C),M_{2^k}(\QU)$) zu erhalten,
genügt es, eine lineare Abbildung $\vartheta:V\to A$ anzugeben mit
\[
\vartheta(v)^2=-Q(v)\,1_A\quad(\forall v\in V).
\]
Dann existiert \emph{eindeutig} ein Algebrahomomorphismus
$\Phi:\Cl(V,Q)\to A$ mit $\Phi\circ\iota=\vartheta$.
Ist $\Phi$ bijektiv, liefert dies einen expliziten Isomorphismus
$\Cl(V,Q)\cong A$ (Bakers Konstruktionen in Kap.~5 arbeiten genau so).
\end{CONC}

\begin{EXA}{Reelle Standardreihe (kleine $n$)}{exa:small-n}
Für $Q=\|\cdot\|^2$:
\[
\Cl_0\cong\R,\qquad
\Cl_1\cong\C,\qquad
\Cl_2\cong\QU,\qquad
\Cl_3\cong\QU\oplus\QU,\qquad
\Cl_4\cong M_2(\QU),
\]
wobei die Abbildungen $V\to A$ explizit gewählt werden können (z.\,B. durch
geeignete Pauli-/Dirac-Blöcke oder Bakers $\vartheta_4:\R^4\to M_2(\QU)$).
\end{EXA}

\begin{REM}{„1-zu-1“-Übersetzung — Merkkasten}{rem:one-to-one}
\begin{itemize}[leftmargin=*]
\item \textbf{Baker:} „Erzeuge $e_i$ mit $e_i^2=-1$ und $e_ie_j=-e_je_i$“\\
\(\Longleftrightarrow\) \textbf{Quotient:} „Setze $I_Q=\langle v\otimes v+Q(v)1\rangle$ und
$\Cl(V,Q)=T(V)/I_Q$; für ONB sind die Relationen die von Baker.“
\item \textbf{Baker:} „Gib ein Modell, indem du konkrete Matrizen $F_i$ mit $F_i^2=-1$ und $F_iF_j=-F_jF_i$ angibst.“\\
\(\Longleftrightarrow\) \textbf{Quotient:} „Gib $\vartheta:V\to A$ mit $\vartheta(v)^2=-Q(v)1$ an —
die Universaleigenschaft liefert den eindeutigen Algebrahomomorphismus.“
\item \textbf{Baker:} „Basis aus sortierten Produkten $e_{i_1}\cdots e_{i_r}$.“\\
\(\Longleftrightarrow\) \textbf{Quotient:} „Filtration $\Rightarrow$ $\mathrm{gr}(\Cl)\cong\Lambda(V)$
$\Rightarrow$ monomische Basis ($2^n$-Dimensionalität).“
\end{itemize}
\end{REM}



\begin{DEF}{Reelle Clifford-Algebra $\Cl_n$}{def-clifford}
Sei $V=\R^n$ mit der Standardform $Q(x)=\|x\|^2=x_1^2+\cdots+x_n^2$.
Die reelle Clifford-Algebra $\Cl_n$ ist die assoziative $\R$-Algebra
mit $1$, erzeugt von Elementen $e_1,\dots,e_n$ und den Relationen
\[
e_i^2=-1,\qquad e_ie_j=-e_je_i\quad(i\neq j).
\]
Als $\R$-Vektorraum besitzt $\Cl_n$ die Basis der geordneten Monome
$e_{i_1}\cdots e_{i_r}$ ($0\le r\le n$, $i_1<\dots<i_r$).
\end{DEF}

\begin{REM}{Universelle Eigenschaft}{universal-property}
Die Einbettung $j_n:\R^n\hookrightarrow \Cl_n$, $x\mapsto j_n(x)$,
erfüllt $j_n(x)^2=-\|x\|^2$.
\emph{Universell:} Für jede $\R$-Algebra $A$ und $\R$-lineares
$f:\R^n\to A$ mit $f(x)^2=-\|x\|^2$ existiert genau ein
Algebrahomomorphismus $F:\Cl_n\to A$ mit $F\circ j_n=f$.
Dies ist die \emph{universelle Eigenschaft} der Clifford-Algebra
(vgl.\ Baker, Thm.~5.4).
\end{REM}

\begin{PROOF}{Skizze}{proof-universal}
Definiere $F(e_{i_1}\cdots e_{i_r})=f(e_{i_1})\cdots f(e_{i_r})$.
Die Relationen garantieren Wohldefiniertheit,
und die Monome $e_{i_1}\cdots e_{i_r}$ bilden eine Basis,
woraus Eindeutigkeit folgt.
\end{PROOF}

\begin{CONC}{Bedeutung der universellen Eigenschaft}{principle-universal}
Um ein konkretes Modell von $\Cl_n$ in einer bekannten Algebra $A$
(z.\,B.\ einer Matrixalgebra) zu erhalten,
reicht es, ein $\R$-lineares $f:\R^n\to A$ mit $f(x)^2=-\|x\|^2$
anzugeben.
Dann induziert die universelle Eigenschaft einen eindeutigen
Algebrahomomorphismus $F:\Cl_n\to A$.
Ist $F$ bijektiv, so erhält man eine explizite Darstellung
\[
\Cl_n\;\cong\;A.
\]
\end{CONC}

\begin{EXA}{Bakers explizite Modelle}{exa-baker-models}
Baker konstruiert lineare Abbildungen
$\vartheta_n:\R^n\to M_{2^k}(\R)$ bzw.\ $M_{2^k}(\C)$ oder $M_{2^k}(\QU)$
mit $\vartheta_n(x)^2=-\|x\|^2$, z.\,B.:
\begin{align*}
\Cl_1&\cong \C,\\
\Cl_2&\cong \QU,\\
\Cl_4&\cong M_2(\QU),
\end{align*}
wobei die letzte Isomorphie durch eine explizite Abbildung
$\vartheta_4:\R^4\to M_2(\QU)$ realisiert wird.
Dies folgt jeweils aus der universellen Eigenschaft
(Baker, Thm.~5.4).
\end{EXA}

\begin{THEO}{Acht-Periodizität}{thm-periodicity}
Für alle $n\ge0$ gilt der Isomorphismus
\[
\Cl_{n+8}\;\cong\;M_{16}(\Cl_n),
\]
wobei $M_k(\Cl_n)$ die Matrixalgebra über $\Cl_n$ bezeichnet.
Dies ergibt die reelle \emph{Bott-Periodizität} auf Algebren-Ebene.
(Baker, Thm.~5.7.)
\end{THEO}

\begin{CONC}{Algebraische Ableitung von Spin- und Pin-Gruppen}{spin-pin}
Innerhalb der Einheitengruppe $\Cl_n^\times$ definieren Produkte
von Einheitsvektoren $u_1\cdots u_k$ die
\emph{Pin-Gruppe} $\operatorname{Pin}(n)$;
die Produkte gerader Länge bilden die \emph{Spin-Gruppe} $\Spin(n)$.
Die Konjugationswirkung
\[
x\;\mapsto\;u\,x\,u^{-1},\qquad u\in\operatorname{Pin}(n),
\]
induziert einen surjektiven Gruppenhomomorphismus
\[
p:\operatorname{Pin}(n)\twoheadrightarrow O(n),\qquad
p|_{\Spin(n)}:\Spin(n)\twoheadrightarrow SO(n),
\]
mit Kern $\{\pm1\}$.
Damit ist $\Spin(n)$ die kanonische zweifache Überlagerung von $SO(n)$
(Baker, Thm.~5.27–5.28).
\end{CONC}

\begin{PROP}{Lie-Algebra $\mathfrak{spin}(n)$}{spin-alg}
Die zugehörige Lie-Algebra innerhalb $\Cl_n$ ist
\[
\mathfrak{spin}(n)
=\mathrm{span}_\R\{\,e_i e_j:\ 1\le i<j\le n\,\},
\]
und das Differential $dp:\mathfrak{spin}(n)\to\mathfrak{so}(n)$
ist ein Isomorphismus.
Die Exponentialkurven $\exp(t\,e_ie_j)\in\Spin(n)$ erzeugen
genau die Planrotationen in der $(i,j)$-Ebene:
$p(\exp(t\,e_ie_j))$ ist die Rotation mit Winkel $2t$.
(Baker, Prop.~5.29–5.30.)
\end{PROP}

\begin{REM}{Zusammenfassung: Algebraischer Zugang}{summary-baker}
\begin{itemize}[leftmargin=*]
\item Die Clifford-Algebra ist ein \emph{Quotient der Tensoralgebra}
  $T(V)$ durch das Ideal
  $\langle v\otimes v+Q(v)\,1\rangle$.
\item Die universelle Eigenschaft erlaubt jede Realisierung in
  einer anderen Algebra durch Angabe von $f(v)^2=-Q(v)$.
\item Bakers Matrizenmodelle (z.\,B.\ $\vartheta_4$) sind genau solche
  Realisierungen und liefern konkrete Isomorphismen
  $\Cl_n\cong M_{k}(\R)$, $M_{k}(\C)$ oder $M_{k}(\QU)$.
\item Die Spin-/Pin-Gruppen entstehen als Untergruppen von
  $\Cl_n^\times$; die Überlagerung $\Spin(n)\to SO(n)$
  folgt rein algebraisch aus der Konjugationswirkung.
\end{itemize}
\end{REM}



\subsection{Vom abstrakten Clifford-Setting zum Matrixmodell (und zurück)}
\label{sec:clifford-matrixbridge}

\begin{MOT}{Motivation: Zwei Welten, eine Struktur}{mot-matrixbridge}
Die Clifford-Algebra $\Cl(V,Q)$ besitzt zwei gleichwertige, aber
komplementäre Realisierungen:
\begin{itemize}[leftmargin=*]
\item[(i)] \emph{Abstrakt:} als universeller Quotient $T(V)/\langle v\otimes v+Q(v)1\rangle$,
  mit klarer funktorieller Bedeutung.
\item[(ii)] \emph{Konkret:} als Unteralgebra oder Untergruppe einer Matrixalgebra
  (Bakers „Matrixgruppen-Ansatz“), die dieselben Relationen realisiert.
\end{itemize}
Ziel ist es, eine natürliche Prozedur zu formulieren, die erlaubt,
zwischen diesen beiden Beschreibungen zu wechseln.
\end{MOT}

\emph{Schritt 1: Vom universellen Objekt zur konkreten Matrixdarstellung}

\begin{DEF}{Startpunkt: Universelle Definition}{universal-start}
Ausgangspunkt ist die universelle Clifford-Algebra
\[
\Cl(V,Q) = T(V)\big/\langle v\otimes v + Q(v)1\rangle,
\]
mit kanonischer Einbettung $\iota:V\hookrightarrow \Cl(V,Q)$
und Relation $\iota(v)^2=-Q(v)1$.
\end{DEF}

\begin{REM}{Ziel: eine Darstellung in $M_k(R)$ finden}{matrix-goal}
Wir möchten ein \emph{konkretes Modell}
\[
\Phi:\Cl(V,Q)\longrightarrow M_k(R)
\]
konstruieren, das ein Isomorphismus (oder wenigstens ein Monomorphismus) ist.
Dazu genügt nach der universellen Eigenschaft, eine lineare Abbildung
\[
f:V\longrightarrow M_k(R)
\]
zu finden, die dieselbe quadratische Relation erfüllt:
\[
f(v)^2 = -\,Q(v)\,I_k \quad (\forall v\in V).
\]
Dann existiert eindeutig ein Algebrahomomorphismus $\Phi$ mit $\Phi\circ\iota=f$.
\end{REM}

\begin{CONC}{Bauprinzip: Von Generatoren zu Matrizen}{prin-gen-to-matrix}
\begin{enumerate}[label=(\alph*), leftmargin=*]
\item Wähle eine Basis $(e_1,\dots,e_n)$ von $V$.
\item Suche Matrizen $F_i\in M_k(R)$ mit
  \[
  F_i^2 = -\,Q(e_i)\,I_k,\qquad F_iF_j = -\,F_jF_i \text{ für } i\neq j.
  \]
\item Definiere $f(e_i):=F_i$ und erweitere $f$ linear.
\item Durch die universelle Eigenschaft ergibt sich
  \[
  \Phi_f:\Cl(V,Q)\longrightarrow M_k(R),\quad
  \Phi_f(e_{i_1}\cdots e_{i_r})=F_{i_1}\cdots F_{i_r}.
  \]
\item Ist $\Phi_f$ injektiv oder surjektiv, erhält man ein konkretes Modell
  (Beispiel: Baker~\cite[Thm.~5.4]{Baker2002}).
\end{enumerate}
\end{CONC}

\begin{EXA}{Beispiel: $\Cl_2$ und die Quaternionen}{exa-quaternions}
Für $V=\R^2$, $Q(x,y)=x^2+y^2$, wähle
\[
F_1=\begin{pmatrix}0&-1\\1&0\end{pmatrix},\quad
F_2=\begin{pmatrix}0&1\\1&0\end{pmatrix}.
\]
Dann $F_i^2=-I$, $F_1F_2=-F_2F_1$, und der induzierte
Homomorphismus $\Phi:\Cl_2\to M_2(\R)$ ist ein Isomorphismus
\[
\Cl_2\;\cong\;\mathbb{H},
\]
wobei $\mathbb{H}$ über die Identifikation
$\begin{smallmatrix}i\mapsto F_1,\\ j\mapsto F_2,\\ k\mapsto F_1F_2\end{smallmatrix}$ repräsentiert wird.
\end{EXA}

\begin{EXA}{Beispiel: Bakers Konstruktion für $\Cl_4$}{exa-baker-4}
Baker definiert
\[
\vartheta_4:\R^4\longrightarrow M_2(\mathbb{H}),\qquad
(x_1,\dots,x_4)\mapsto
\begin{bmatrix}
x_1 i+x_2 j+x_3 k & x_4 k\\
x_4 k & x_1 i+x_2 j - x_3 k
\end{bmatrix}.
\]
Dann gilt $\vartheta_4(x)^2=-\|x\|^2I$, also induziert die universelle Eigenschaft
\[
\Phi:\Cl_4\longrightarrow M_2(\mathbb{H}),
\]
und $\Phi$ ist ein Isomorphismus.
So werden die abstrakten Relationen konkret in Matrixmultiplikation kodiert.
\end{EXA}

\begin{REM}{Interpretation: von der Theorie zur Praxis}{rem-theory-to-practice}
Diese Konstruktion entspricht exakt der Funktorialität des freien Objekts:
$T(V)$ repräsentiert alle bilinearen Multiplikationen, und
die Relation $v^2=-Q(v)$ schneidet den universellen Raum auf
jene Repräsentationen zu, die die Metrik respektieren.
Das Einsetzen konkreter Matrizen $(F_i)$ implementiert
eine konkrete Realisierung dieses universellen Objekts
als Matrixalgebra.
\end{REM}

\emph{Schritt 2: Vom Matrixmodell zurück zur abstrakten Algebra}

\begin{DEF}{Rückrichtung: Abstraktion}{def-backtoabs}
Gegeben sei eine Familie von Matrizen $(F_i)\subset M_k(R)$ mit
\[
F_i^2=-1,\qquad F_iF_j=-F_jF_i.
\]
Diese erzeugen eine Unteralgebra $A\subset M_k(R)$.
Dann gibt es nach der universellen Eigenschaft \emph{einen eindeutigen}
Surjektionshomomorphismus
\[
\Psi:\Cl_n\longrightarrow A,\qquad \Psi(e_i)=F_i.
\]
Das Bild $A$ ist \emph{isomorph zur kanonischen Clifford-Algebra} $\Cl_n/\ker\Psi$.
\end{DEF}

\begin{PROP}{Abstraktion als kanonischer Vergiss-Funktor}{prop-forget}
Die Zuordnung
\[
A=(F_1,\dots,F_n)\subset M_k(R)\quad\longmapsto\quad (\Cl_n,\Psi)
\]
ist ein Funktor vom \emph{Matrix-Modellraum} in die Kategorie
$\Alg_\R$ der reellen Algebren, der die universelle Eigenschaft respektiert:
Morphismen von Matrizen, die Relationen erhalten, induzieren
Algebrahomomorphismen zwischen den zugehörigen Clifford-Algebren.
\end{PROP}

\begin{REM}{Praktische Bedeutung}{rem-practical}
Das bedeutet: Wenn man eine konkrete Darstellung (z.\,B.\ in $M_2(\mathbb{H})$)
berechnet, kann man stets \emph{abstrakt} zurücktreten und sagen:
„Das ist \(\Cl_4\) in einem Modell.“  
Wechselt man später zu einer anderen Realisierung (etwa in $M_4(\C)$),
so ist das nur ein anderer Repräsentant desselben universellen Objekts.
\end{REM}

\emph{Schritt 3: Der Wechsel der Modelle als Funktor (abstrakt $\leftrightarrow$ konkret)}

\begin{CONC}{Zusammenfassung: bidirektionale Prozedur}{conc-procedure}
\begin{center}
\begin{tikzcd}[row sep=huge, column sep=large]
(V,Q) \arrow[r, mapsto, "\text{freies Objekt}"] & 
\Cl(V,Q)=T(V)/\langle v^2+Q(v)1\rangle
\arrow[r, dashed, "\exists!\,f\text{ mit }f(v)^2=-Q(v)I"] &
A\subseteq M_k(R)
\end{tikzcd}
\end{center}
\begin{enumerate}[leftmargin=*]
\item \textbf{Abstrakt $\to$ Konkret:}\\
  Beginne mit $V,Q$, wähle eine Basis und eine konkrete lineare Abbildung
  $f:V\to M_k(R)$ mit $f(v)^2=-Q(v)I$.  
  Dann liefert die universelle Eigenschaft den eindeutigen Homomorphismus
  $\Phi:\Cl(V,Q)\to M_k(R)$, der die abstrakte Algebra realisiert.

\item \textbf{Konkret $\to$ Abstrakt:}\\
  Gegeben eine Familie $(F_i)$ von Matrizen mit Bakers Relationen,
  definiere $\Psi:\Cl_n\to M_k(R)$ durch $\Psi(e_i)=F_i$.  
  Das Bild ist eine Repräsentation (Matrixmodell) des universellen Objekts.

\item \textbf{Wechsel der Modelle:}\\
  Jede zwei Realisierungen $(A,f)$ und $(A',f')$ sind durch eindeutige
  Morphismen von $\Cl(V,Q)$ verbunden:
  \[
  \Hom_{\Alg_\R}(\Cl(V,Q),A)\;\cong\;
  \{f:V\to A\mid f(v)^2=-Q(v)\}.
  \]
  Der Wechsel zwischen Modellen entspricht also einfach dem Wechsel
  der gewählten linearen Einbettung $f$.
\end{enumerate}
\end{CONC}

\begin{REM}{Fazit: Natürlicher Wechsel zwischen beiden Welten}{rem-summary}
Nach Verinnerlichung dieser Prozedur kann man intuitiv
„zwischen den Welten“ wechseln:
\begin{itemize}[leftmargin=*]
\item Die abstrakte Sicht zeigt, \emph{welche Struktur} man modelliert
  (universelle Relationen, Funktorialität, Spin/Pin).
\item Die Matrixsicht liefert \emph{wie} diese Struktur in konkreten Rechnungen
  (z.\,B.\ in $\mathrm{M}_2(\mathbb{H})$ oder $\mathrm{M}_4(\C)$) implementiert ist.
\item Beides beschreibt dasselbe Objekt — die Clifford-Algebra
  als universelles Trägermodell einer bilinearen Form.
\end{itemize}
Damit wird das Springen zwischen den Settings
\[
\textit{„Abstrakt – Funktoriell – Geometrisch“}\quad\longleftrightarrow\quad
\textit{„Konkret – Rechnerisch – Matrizial“}
\]
nicht zu einem Wechsel der Theorie, sondern nur zu einem Wechsel der Darstellung.
\end{REM}


\pagebreak


\section{Kohomologiealgebren und das Cup-Produkt}
\label{sec:cohomology-algebra}

\begin{MOT}{Zielbild}{cohomology-goal}
Wir wollen präzisieren, dass singuläre Kohomologie \(H^\ast(X;R)\) eines Raums \(X\) eine \emph{gradierte, superkommutative \(R\)-Algebra} ist, und die Ringstruktur \(\smile\) \emph{kanonisch} aus Kettenniveau-Daten (\(\Delta\) bzw.\ Alexander--Whitney) entsteht. Parallel betonen wir Funktorialität, Kreuz-/Cap-Produkt und den monoidalkategorischen Blick.
\end{MOT}

\subsection{Ketten- und Kokettenkomplex; Differential graded algebra (DGA)}

\begin{DEF}{Singuläre Ketten und Koketten}{chains-cochains}
Sei \(R\) ein kommutativer Ring. Der singuläre \(n\)-Kettenmodul von \(X\) ist
\[
C_n(X;R):=R\big\langle \sigma:\Delta^n\to X\ \text{singulär}\big\rangle,\qquad
\partial=\sum_{i=0}^{n}(-1)^i\,d_i: C_n\to C_{n-1}.
\]
Der Kohomologie-\emph{Koketten}komplex ist der \(R\)-duale
\[
C^n(X;R):=\Hom_R\!\big(C_n(X;R),R\big),
\]
mit Kohomologie-Differential
\[
\delta:C^n\to C^{n+1},\qquad (\delta\varphi)(c)=\varphi(\partial c).
\]
\end{DEF}

\begin{DEF}{Alexander--Whitney-Diagonale}{AW}
Die Alexander--Whitney-Abbildung ist ein natürlicher Kettenmorphismus
\[
\mathrm{AW}:C_\ast(X;R)\longrightarrow C_\ast(X;R)\otimes_R C_\ast(X;R)
\]
definiert auf einer singulären Simplex \(\sigma:\Delta^n\to X\) durch
\[
\mathrm{AW}(\sigma)=\sum_{p+q=n}\,
\sigma\!\circ\! [v_0,\dots,v_p]\ \otimes\ \sigma\!\circ\! [v_p,\dots,v_n],
\]
wobei \([v_i,\dots,v_j]\) die kanonisch einbettende Seiten-Simplex bezeichnet. Man hat \((\partial\otimes 1+1\otimes\partial)\circ \mathrm{AW}=\mathrm{AW}\circ\partial\).
\end{DEF}

\begin{DEF}{Cup-Produkt auf Koketten}{cup-on-cochains}
Für \(\varphi\in C^p(X;R),\ \psi\in C^q(X;R)\) definiert
\[
\varphi\smile\psi\ :=\ (\varphi\otimes\psi)\circ \mathrm{AW}\ \in C^{p+q}(X;R).
\]
Damit wird \((C^\ast(X;R),\delta,\smile)\) zu einer \emph{(nicht streng kommutativen) DGA}:
\[
\delta(\varphi\smile\psi)=\delta\varphi\smile\psi+(-1)^p\,\varphi\smile\delta\psi,\qquad
(\varphi\smile\psi)\smile\chi=\varphi\smile(\psi\smile\chi).
\]
Das Einheits\-element ist \(1\in C^0(X;R)\), gegeben durch die konstante Funktion \(1\) auf \(0\)-Ketten.
\end{DEF}

\begin{REM}{Superkommutativität bis auf Homotopie}{graded-comm-up-to-homotopy}
Auf Kettenniveau ist \(\smile\) i.\,A.\ nicht (super-)kommutativ, aber es existiert eine explizite Kettenhomotopie \(H\) mit
\[
\varphi\smile\psi-(-1)^{pq}\psi\smile\varphi=\delta H(\varphi,\psi)+H(\delta\varphi,\psi)+(-1)^p H(\varphi,\delta\psi).
\]
Daraus folgt, dass in der Kohomologie \(H^\ast(X;R)\) die \emph{gradierte Kommutativität} gilt (siehe Proposition \ref{prop:cup-properties}).
\end{REM}

\subsection{Ringstruktur auf der Kohomologie und Funktorialität}

\begin{DEF}{Kohomologiealgebra}{cohomology-algebra}
Die singuläre Kohomologie \(H^\ast(X;R)=\bigoplus_{n\ge 0}H^n(X;R)\) erhält über Repräsentanten die Multiplikation
\[
[\alpha]\cdot[\beta]\ :=\ [\alpha\smile\beta]\ \in H^{p+q}(X;R),
\quad \alpha\in Z^p,\ \beta\in Z^q.
\]
Dies ist wohldefiniert aufgrund der DGA-Regel \(\delta(\alpha\smile\beta)=0\) für Zyklen und der Homotopie-Argumente oben.
\end{DEF}

\begin{PROP}{Eigenschaften des Cup-Produkts}{prop:cup-properties}
Für \(\alpha\in H^p(X;R)\), \(\beta\in H^q(X;R)\), \(\gamma\in H^r(X;R)\) gilt:
\begin{romanenum}
\item \textbf{Bilinearität:} \([\alpha_1+\alpha_2]\smile\beta=\alpha_1\smile\beta+\alpha_2\smile\beta\).
\item \textbf{Assoziativität:} \((\alpha\smile\beta)\smile\gamma=\alpha\smile(\beta\smile\gamma)\).
\item \textbf{Einheit:} \(1\in H^0(X;R)\) ist die multiplikative Einheit.
\item \textbf{Gradierte Kommutativität:} \(\alpha\smile\beta=(-1)^{pq}\beta\smile\alpha\).
\item \textbf{Functorialität:} Für stetig \(f:X\to Y\) ist \(f^\ast:H^\ast(Y;R)\to H^\ast(X;R)\) ein \emph{Grad-\(0\)-Algebrahomomorphismus}:
\(f^\ast(\alpha\smile\beta)=f^\ast\alpha\smile f^\ast\beta\).
\end{romanenum}
\end{PROP}

\begin{REM}{Kategorischer Blick}{cat-view}
Damit ist \(H^\ast(-;R):\mathbf{Top}^{\mathrm{op}}\to \Alg_R^{\mathrm{gr,\,supercomm}}\) ein kontravarianter Funktor in die Kategorie der graduiert superkommutativen \(R\)-Algebren. Auf Kettenniveau ist \(C^\ast(-;R)\) ein Funktor in die Kategorie der DGAs.
\end{REM}

\subsection{Nachweis der Algebraeigenschaft des Kohomologierings}

\begin{DEF}{Singuläre Cochains, Differential, Cup-Produkt}{cochains-cup}
Sei $X$ ein topologischer Raum und $R$ ein kommutativer Ring. 
Wir setzen $C_n(X;R)$ für die singulären $n$-Ketten und $C^n(X;R):=\Hom_R(C_n(X;R),R)$.
Das Coboundary ist $\delta:C^n\to C^{n+1}$, $\;\delta\alpha:=\alpha\circ\partial$.

Die \emph{Alexander–Whitney}–Diagonale $\mathrm{AW}:C_{p+q}(X)\to C_p(X)\otimes C_q(X)$ ist koassoziativ und natürlich. 
Für $\alpha\in C^p,\ \beta\in C^q$ definieren wir das Cup-Produkt kochainweise durch
\[
\alpha\smile\beta\;:=\;(\alpha\otimes\beta)\circ \mathrm{AW}\;\in C^{p+q}(X;R).
\]
\end{DEF}

\begin{LEM}{Leibniz-Regel auf Cochain-Ebene}{leibniz}
Für $\alpha\in C^p,\ \beta\in C^q$ gilt
\[
\delta(\alpha\smile\beta)\;=\;(\delta\alpha)\smile\beta + (-1)^p\,\alpha\smile(\delta\beta).
\]
\end{LEM}

\begin{PROOF}{zu Lemma \ref{leibniz}}{proof-leibniz}
Schreibe $\alpha\smile\beta=(\alpha\otimes\beta)\circ \mathrm{AW}$. 
Da $\partial$ ein Derivationsgesetz zu $\mathrm{AW}$ erfüllt (Dual von Koalgebra-Axiomen), erhalten wir
\[
\delta(\alpha\smile\beta)
=(\alpha\otimes\beta)\circ (\partial\otimes\id+\;(-1)^p\,\id\otimes\partial)\circ\mathrm{AW}
=(\delta\alpha)\smile\beta + (-1)^p\,\alpha\smile(\delta\beta).
\qedhere
\]
\end{PROOF}

\begin{PROP}{Wohldefiniertheit auf Kohomologie}{well-defined}
Ist $\alpha\sim \alpha+\delta\xi$ und $\beta\sim \beta+\delta\eta$ kohomolog, so gilt
\[
[\alpha\smile\beta]=[(\alpha+\delta\xi)\smile(\beta+\delta\eta)]\quad\text{in }H^{p+q}(X;R).
\]
\end{PROP}

\begin{PROOF}{zu Proposition \ref{well-defined}}{proof-well-defined}
Durch Bilinearität und Lemma~\ref{leibniz}:
\[
(\alpha+\delta\xi)\smile\beta-\alpha\smile\beta
=\delta\xi\smile\beta
=\delta(\xi\smile\beta)-(-1)^{p-1}\xi\smile\delta\beta.
\]
Für Cocycle $\beta$ ist der Unterschied ein Coboundary. Analog für Variation von $\beta$.
\end{PROOF}

\begin{PROP}{Assoziativität und Einheit}{assoc-unit}
Das Cup-Produkt ist assoziativ auf Cochains und damit auf Kohomologie:
\[
(\alpha\smile\beta)\smile\gamma=\alpha\smile(\beta\smile\gamma).
\]
Ferner ist die Klasse $[1]\in H^0(X;R)$ (der konstante $0$-Cochain $1$) eine multiplikative Einheit.
\end{PROP}

\begin{PROOF}{zu Proposition \ref{assoc-unit}}{proof-assoc-unit}
Koassoziativität von $\mathrm{AW}$ liefert
\[
(\alpha\smile\beta)\smile\gamma
=(\alpha\otimes\beta\otimes\gamma)\circ(\mathrm{AW}\otimes\id)\circ\mathrm{AW}
=(\alpha\otimes\beta\otimes\gamma)\circ(\id\otimes\mathrm{AW})\circ\mathrm{AW}
=\alpha\smile(\beta\smile\gamma).
\]
Für die Einheit beachte: $\mathrm{AW}$ hat $R\cong C_0$ als Kounität; dual ist der konstante Cochain $1$ eine Einheit, d.\,h.\ $1\smile\alpha=\alpha=\alpha\smile 1$.
\end{PROOF}

\begin{PROP}{Gradierte Kommutativität auf Kohomologie}{graded-comm}
Für Kohomologieklassen $[\alpha]\in H^p,\ [\beta]\in H^q$ gilt
\[
[\alpha]\smile[\beta]\;=\;(-1)^{pq}\,[\beta]\smile[\alpha].
\]
\end{PROP}

\begin{PROOF}{zu Proposition \ref{graded-comm}}{proof-graded-comm}
Auf Cochain-Ebene ist $\smile$ nicht kommutativ. 
Es existiert jedoch die natürliche \emph{cup-eins}–Operation $\smallsmile_1:C^p\times C^q\to C^{p+q-1}$ (Steenrod), 
so dass (für homogene $\alpha,\beta$) die Homotopiegleichung gilt:
\[
\delta(\alpha\smallsmile_1\beta)
=(-1)^p\,\alpha\smile\beta - (-1)^{pq}\,\beta\smile\alpha
-(-1)^p(\delta\alpha)\smallsmile_1\beta - \alpha\smallsmile_1(\delta\beta).
\]
Sind $\alpha,\beta$ Cocycle, vereinfacht sich dies zu
\[
\alpha\smile\beta - (-1)^{pq}\,\beta\smile\alpha
=\delta\big((-1)^p\,\alpha\smallsmile_1\beta\big),
\]
also ist der Unterschied ein Coboundary; die behauptete Gleichheit folgt in $H^\ast(X;R)$.
\end{PROOF}

\begin{PROP}{Functorialität}{functoriality}
Für $f:X\to Y$ ist $f^*:H^\ast(Y;R)\to H^\ast(X;R)$ ein Ringhomomorphismus graduierter Ringe:
\[
f^*(\alpha\smile\beta)=f^*\alpha\smile f^*\beta,\qquad f^*(1)=1.
\]
\end{PROP}

\begin{PROOF}{zu Proposition \ref{functoriality}}{proof-functoriality}
Naturality von $\mathrm{AW}$ liefert $(f_*\otimes f_*)\circ \mathrm{AW}_X=\mathrm{AW}_Y\circ f_*$.
Dual folgt $(f^*\alpha)\smile(f^*\beta)=f^*(\alpha\smile\beta)$. Die Einheit ist offensichtlich.
\end{PROOF}

\begin{CONC}{Algebra-Struktur von $H^\ast(X;R)$}{cohom-alg}
Mit den obigen Lemmata/Propositionen ist $H^\ast(X;R)$ eine \emph{graduierte, superkommutative $R$-Algebra} 
mit Einheit $[1]\in H^0(X;R)$.
\end{CONC}


\subsection{Kreuzprodukt, Diagonale und Künneth-Bezug}

\begin{DEF}{Eilenberg--Zilber und Kreuzprodukt}{cross}
Der Eilenberg--Zilber-Isomorphismus \( \mathrm{EZ}:C_\ast(X\times Y)\xrightarrow{\simeq} C_\ast(X)\otimes C_\ast(Y)\) ist ein natürlicher Kettenäquivalenz. Dual erhält man ein (auf Kohomologie wohldefiniertes) \emph{Kreuzprodukt}
\[
\times:\ H^p(X;R)\otimes H^q(Y;R)\longrightarrow H^{p+q}(X\times Y;R).
\]
\end{DEF}

\begin{PROP}{Cup-Produkt als Pullback des Kreuzprodukts}{diag}
Schreibe \(\Delta:X\to X\times X\) für die Diagonale. Dann gilt
\[
\alpha\smile\beta=\Delta^\ast(\alpha\times\beta)\ \in H^{p+q}(X;R).
\]
\end{PROP}

\begin{REM}{Künneth-Heuristik}{kuenneth}
Unter geeigneten Voraussetzungen (z.\,B.\ flache Koeffizienten) liefert der Künneth-Satz eine Beschreibung von \(H^\ast(X\times Y;R)\) via \(H^\ast(X;R)\) und \(H^\ast(Y;R)\), kompatibel mit \(\times\).
\end{REM}

\subsection{Cap-Produkt und Modulstruktur der Homologie}

\begin{DEF}{Cap-Produkt}{cap}
Das Cap-Produkt
\[
\smallsmile:\ H^p(X;R)\otimes H_q(X;R)\longrightarrow H_{q-p}(X;R)
\]
wird auf Kettenniveau via \(\mathrm{AW}\) (bzw.\ seine Adjunktion) definiert. Es macht \(H_\ast(X;R)\) zu einem \emph{linksmodul} über der Kohomologiealgebra \(H^\ast(X;R)\):
\[
(\alpha\smile\beta)\smallsmile c=\alpha\smallsmile(\beta\smallsmile c),\qquad
1\smallsmile c=c.
\]
\end{DEF}

\begin{REM}{Poincaré-Dualität (orientierter, kompakter, glatter Fall)}{PD}
Ist \(M^n\) eine orientierte, kompakte, zusammenhängende \(n\)-dimensionale Mannigfaltigkeit und \(R\) ein Körper, so induziert die Fundamentalklasse \([M]\in H_n(M;R)\) einen Isomorphismus
\[
\mathrm{PD}:H^k(M;R)\xrightarrow{\ \cong\ } H_{n-k}(M;R),\qquad \alpha\mapsto \alpha\smallsmile [M],
\]
kompatibel mit der Modulstruktur.
\end{REM}

\subsection{Beispiele und Rechenregeln}

\begin{EXA}{Sphären und Produkte}{spheres}
\begin{itemize}[leftmargin=*]
\item \(H^\ast(S^n;R)\cong R[\alpha]/(\alpha^2)\) mit \(\deg\alpha=n\).
\item \(H^\ast(S^p\times S^q;R)\cong R[\alpha,\beta]/(\alpha^2,\beta^2)\) mit \(\deg\alpha=p,\ \deg\beta=q\), und \(\alpha\beta=(-1)^{pq}\beta\alpha\).
\end{itemize}
\end{EXA}

\begin{REM}{Fortgeschritten: \(E_\infty\)-Struktur}{einfty}
Der Kohomologie-DGA \(C^\ast(X;R)\) trägt eine kanonische \(E_\infty\)-Algebra-Struktur (höhere Homotopien für Kommutativität/Assoziativität). Auf \(H^\ast\) „stürzt“ dies zu einer streng superkommutativen Algebra ab; die \(E_\infty\)-Struktur wird in Operaden- und Homotopietheorie wesentlich.
\end{REM}

\subsection{Zusammenfassung (Funktorielle Ringstruktur)}
\begin{CONC}{Cohomology as a graded-commutative ring functor}{summary}
Die Daten \((C_\ast,\partial)\), \(\mathrm{AW}\), \(\mathrm{EZ}\) und \(\Delta\) liefern kanonisch:
\[
H^\ast(-;R):\ \mathbf{Top}^{\mathrm{op}}\longrightarrow \Alg_R^{\mathrm{gr,\,supercomm}},
\]
mit
\begin{enumerate}[leftmargin=*]
\item Ringstruktur durch \(\smile\) (Definition über \(\mathrm{AW}\));
\item Produktverträglichkeit zu Abbildungen (\(f^\ast\) ist Algebrahom.);
\item Kreuzprodukt \(\times\) und Relation \(\smile=\Delta^\ast\circ\times\);
\item Cap-Produkt, das \(H_\ast\) zum Modul über \(H^\ast\) macht.
\end{enumerate}
\end{CONC}



\subsection{Formelles Vorgehen: Kohomologieringe als Quotienten freier Algebren}
\label{sec:cohom-ring-quotients}

\begin{MOT}{Idee}{cohom-quotient-idea}
Kohomologieringe $H^\ast(X;R)$ sind graduierte, superkommutative $R$-Algebren.
Wie in der abstrakten Algebra werden sie daher als
\emph{freie (graduierte) Algebren modulo Relationen} beschrieben:
\[
H^\ast(X;R)\;\cong\;R[\alpha_1,\dots,\alpha_r]/I,
\]
wobei die $\alpha_i$ topologisch definierte Erzeuger sind
(z.\,B. Fundamentalklassen, Chern-Klassen, harmonische Differentialformen),
und $I$ ein homogenes Ideal von Relationen
(aus Cup-Produkten) bezeichnet.
\end{MOT}

\begin{DEF}{Freie graduierte kommutative Algebra}{free-graded-comm}
Für homogene Variablen $\alpha_i$ mit Graden $|\alpha_i|=d_i$
sei
\[
R[\alpha_1,\dots,\alpha_r]
:= S_R(V_{\bar 0})\otimes_R \Lambda_R(V_{\bar 1}),
\]
wobei $V_{\bar 0}$ (bzw. $V_{\bar 1}$)
der $R$-Modul der Generatoren geraden (bzw.\ ungeraden) Parität ist.
Diese Algebra ist die \emph{freie superkommutative, graduierte $R$-Algebra}
auf den Generatoren $\{\alpha_i\}$:
\[
\alpha_i\alpha_j = (-1)^{|\alpha_i||\alpha_j|}\alpha_j\alpha_i.
\]
\end{DEF}

\begin{DEF}{Kohomologiering als Quotient}{cohom-ring-as-quotient}
Für einen Raum $X$ mit endlich vielen Kohomologieerzeugern
wählt man eine Menge homogener Generatoren
$\alpha_1,\dots,\alpha_r\in H^{>0}(X;R)$, so dass
\[
H^\ast(X;R)
\;\cong\;
R[\alpha_1,\dots,\alpha_r]\big/ I_X,
\]
wobei das Ideal $I_X$ alle \emph{Polynomrelationen zwischen Cup-Produkten}
der Generatoren beschreibt:
\[
I_X
=\big\langle
P(\alpha_1,\dots,\alpha_r)\in R[\alpha_1,\dots,\alpha_r]\mid
P(\alpha_1,\dots,\alpha_r)=0\text{ in }H^\ast(X;R)
\big\rangle.
\]
\end{DEF}

\begin{REM}{Topologischer Ursprung der Relationen}{relations}
Die Relationen im Ideal $I_X$ spiegeln geometrische Eigenschaften von $X$ wider:
\begin{itemize}[leftmargin=*]
\item Verschwinden bestimmter Cup-Produkte, z.\,B.\ $\alpha^2=0$,
      folgt aus $H^{2|\alpha|}(X;R)=0$.
\item Potenzrelationen (z.\,B.\ $x^{n+1}=0$ für $\mathbb{CP}^n$)
      entstehen durch Dimensionsbeschränkung der Kohomologie.
\item Lineare Relationen (z.\,B.\ $e_1+e_2+\cdots+e_k=0$)
      treten in Faserbündeln und Klassifizierungsräumen auf.
\end{itemize}
Somit ist das Ideal $I_X$ die algebraische Kodierung
aller topologisch bedingten Multiplikationsrelationen.
\end{REM}

\begin{EXA}{Beispiele}{cohom-ring-quotients}
\begin{enumerate}[label=(\alph*), leftmargin=*]
\item \textbf{Sphäre $S^n$.}
\[
H^\ast(S^n;R)\;\cong\;R[\alpha]/(\alpha^2),
\quad |\alpha|=n.
\]
Nur zwei nichttriviale Grade ($0$ und $n$);
das Produkt $\alpha^2$ liegt in $H^{2n}=0$.

\item \textbf{Produktsphäre $S^p\times S^q$.}
\[
H^\ast(S^p\times S^q;R)
\;\cong\;
R[\alpha,\beta]/(\alpha^2,\beta^2),
\quad |\alpha|=p,\ |\beta|=q,
\quad \alpha\beta=(-1)^{pq}\beta\alpha.
\]
Künneth-Formel und diagonales Cup-Produkt.

\item \textbf{Komplex projektiver Raum.}
\[
H^\ast(\mathbb{CP}^n;R)
\;\cong\;
R[x]/(x^{n+1}),
\quad |x|=2.
\]
Das Cup-Produkt $x^{n+1}$ verschwindet, da $H^{2n+2}=0$.
\end{enumerate}
\end{EXA}

\begin{REM}{Konzeptuelle Analogie zu freien Algebren}{analogy}
Diese Darstellungen sind direkt analog zu den
\emph{freien Algebren} aus Abschnitt~\ref{sec:tensor-algebra}:
\[
H^\ast(X;R)
\simeq
T_R(V_X)\big/ \langle \text{Relations from }X\rangle,
\]
wobei $V_X$ der durch die kohomologischen Erzeuger
aufgespannte $R$-Modul ist.
So entstehen Kohomologieringe formal
als \emph{Quotienten freier, superkommutativer Algebren}
durch das von topologischen Relationen erzeugte Ideal.
\end{REM}

\begin{CONC}{Zusammenhang Algebra–Topologie}{cohom-ring-conclusion}
Das Verfahren
\[
\text{(Topologischer Raum $X$)}\quad\longmapsto\quad
H^\ast(X;R)\cong R[\alpha_1,\dots,\alpha_r]/I_X
\]
verwandelt homologische Daten in eine algebraische Signatur.
Die Generatoren repräsentieren topologische Basisformen,
die Relationen kodieren Geometrie und Dimension,
und die resultierende Polynomquotient-Algebra
ist ein funktorielles, homotopieinvariantes Invariant von $X$.
\end{CONC}

\subsection{Der Kohomologiering von $S^p\times S^q$ als Polynomquotient}

\begin{MOT}{Ziel und Rahmenbedingungen}{SpSq-aim}
Wir bestimmen die graduierte Ringstruktur von $H^\ast(S^p\times S^q;R)$. 
Als Koeffizientenring sei $R$ ein kommutativer Ring mit $1$. 
Da $H_\ast(S^n;\mathbb{Z})$ frei ist, verschwinden in der Künneth-Formel alle $\mathrm{Tor}$-Terme; 
insbesondere gilt (funktoriell in $R$)
\[
H^\ast(S^p\times S^q;R)\;\cong\;H^\ast(S^p;R)\,\widetilde{\otimes}\,H^\ast(S^q;R),
\]
wobei das Tilde-Tensor \(\widetilde{\otimes}\) die \emph{Ring\-tensor\-struktur} mit Zeichenregel (Eilenberg–Zilber/Alexander–Whitney) bezeichnet, siehe \ref{rem:ringtensor}.
\end{MOT}

\begin{DEF}{Erzeuger durch Projektionen}{gens}
Bezeichne die Projektionen mit \(\pi_1:S^p\times S^q\to S^p\), \(\pi_2:S^p\times S^q\to S^q\).
Seien \(\gamma_p\in H^p(S^p;R)\), \(\gamma_q\in H^q(S^q;R)\) die \emph{Fundamentalklassen in Kohomologie} 
(unter Wahl einer Orientierung; über \(R\) via universeller Koeffizienten: \(H^n(S^n;R)\cong R\)). 
Definiere
\[
\alpha := \pi_1^\ast(\gamma_p)\in H^p(S^p\times S^q;R),
\qquad
\beta  := \pi_2^\ast(\gamma_q)\in H^q(S^p\times S^q;R).
\]
\end{DEF}

\begin{REM}{Ko\-algebraischer Ursprung der Multiplikation}{cross-cup}
Die Ringstruktur auf \(H^\ast(S^p\times S^q;R)\) entsteht aus dem \emph{Kreuzprodukt} 
\(\times: H^i(S^p;R)\otimes H^j(S^q;R)\to H^{i+j}(S^p\times S^q;R)\)
und der Diagonale \(\Delta:S^p\times S^q\to(S^p\times S^q)\times(S^p\times S^q)\) via
\[
u\smile v \;=\; \Delta^\ast\big((u_1\times u_2)\big)\quad\text{(Eilenberg–Zilber/AW)}.
\]
Für \(\alpha,\beta\) gilt konkret \(\alpha=\pi_1^\ast(\gamma_p)=\gamma_p\times 1\) und \(\beta=1\times\gamma_q\).
\end{REM}

\begin{REM}{Gradierte Ringstruktur auf dem Tensorprodukt}{ringtensor}
\label{rem:ringtensor}
Auf dem Tensorprodukt zweier \emph{graduiert superkommutativer} Ringe \(A^\ast,B^\ast\) ist die Multiplikation
\[
(a\otimes b)\cdot(a'\otimes b') \;=\; (-1)^{|b|\cdot|a'|}\,(a\smile a')\otimes(b\smile b'),
\]
für homogene \(a,a',b,b'\).
Mit \(A^\ast=H^\ast(S^p;R)\), \(B^\ast=H^\ast(S^q;R)\) erhalten wir so die Ringstruktur auf 
\(H^\ast(S^p\times S^q;R)\cong A^\ast\widetilde{\otimes}B^\ast\).
\end{REM}

\begin{PROP}{Additive Struktur (Künneth)}{additive}
Additiv gilt
\[
H^\ast(S^p\times S^q;R)\;\cong\;
R\cdot 1 \;\oplus\; R\cdot \alpha \;\oplus\; R\cdot \beta \;\oplus\; R\cdot (\alpha\smile\beta),
\]
wobei \(\deg(1)=0\), \(\deg(\alpha)=p\), \(\deg(\beta)=q\), \(\deg(\alpha\beta)=p+q\).
\end{PROP}

\begin{PROOF}{von \ref{additive}}{proof-additive}
Für Sphären gilt \(H^\ast(S^n;R)\cong R\cdot 1\oplus R\cdot \gamma_n\) und sonst \(0\).
Das Künneth-Theorem (ohne \(\mathrm{Tor}\)) liefert
\[
H^\ast(S^p\times S^q;R)\cong
\big(H^\ast(S^p;R)\otimes H^\ast(S^q;R)\big),
\]
mit homogener Basis \(\{1\otimes 1,\,\gamma_p\otimes 1,\,1\otimes\gamma_q,\,\gamma_p\otimes\gamma_q\}\).
Über \(\pi_i^\ast\) identifizieren wir 
\(1\otimes 1\mapsto 1\), \(\gamma_p\otimes 1\mapsto \alpha\), \(1\otimes\gamma_q\mapsto \beta\),
\(\gamma_p\otimes\gamma_q\mapsto \alpha\beta\).
\end{PROOF}

\begin{LEM}{Quadrate verschwinden}{squares-zero}
\(\alpha^2=0\) und \(\beta^2=0\).
\end{LEM}

\begin{PROOF}{}{proof-squares}
\(\alpha\in H^p\), also \(\alpha^2\in H^{2p}(S^p\times S^q;R)\).
Da \(H^{2p}(S^p;R)=0\) und \(H^{>q}(S^q;R)=0\), 
folgt aus dem Künneth-Zerfall, dass \(H^{2p}(S^p\times S^q;R)=0\) (sofern \(2p\notin\{0,p,q,p+q\}\), was bei \(p>0\) stets gilt).
Somit \(\alpha^2=0\); analog \(\beta^2=0\).
(Alternativ: \(\alpha=\gamma_p\times 1\), also \(\alpha^2=(\gamma_p\smile \gamma_p)\times 1=0\times 1=0\), 
weil \(\gamma_p\smile\gamma_p=0\) in \(H^{2p}(S^p;R)\).)
\end{PROOF}

\begin{LEM}{Gradierte Kommutativität}{graded-comm}
\(\alpha\beta = (-1)^{pq}\,\beta\alpha\).
\end{LEM}

\begin{PROOF}{}{proof-graded-comm}
Allgemein ist Kohomologie graduiert superkommutativ:
\(\xi\smile\eta=(-1)^{|\xi||\eta|}\eta\smile\xi\) für homogene \(\xi,\eta\).
Setze \(\xi=\alpha\) (Grad \(p\)), \(\eta=\beta\) (Grad \(q\)).
\end{PROOF}

\begin{LEM}{Nichtverschwindung des Kreuzterms}{nonzero}
\(\alpha\beta\neq 0\) in \(H^{p+q}(S^p\times S^q;R)\cong R\).
\end{LEM}

\begin{PROOF}{}{proof-nonzero}
Unter der Künneth-Isomorphie entspricht \(\alpha\beta\) dem Tensor \(\gamma_p\otimes\gamma_q\).
Da \(H^p(S^p;R)\cong R\cdot\gamma_p\), \(H^q(S^q;R)\cong R\cdot\gamma_q\),
ist \(\gamma_p\otimes\gamma_q\) ein Erzeuger von \(H^{p+q}(S^p\times S^q;R)\cong R\).
Damit ist \(\alpha\beta\) genau die Orientierungsklasse (bis Einheit) und daher \(\neq 0\) (vorausgesetzt \(1_R\neq 0\)).
\end{PROOF}

\begin{PROP}{Ringpräsentation}{presentation}
Mit \(|\alpha|=p\), \(|\beta|=q\) gilt ein kanonischer Ringisomorphismus
\[
\Phi:\; R[\alpha,\beta]/(\alpha^2,\beta^2)\;\xrightarrow{\ \cong\ }\; H^\ast(S^p\times S^q;R),
\]
der die Variablen auf die oben gewählten Klassen abbildet.
\end{PROP}

\begin{PROOF}{}{proof-presentation}
\emph{(1) Surjektivität:} Nach \ref{additive} wird \(H^\ast\) durch \(1,\alpha,\beta,\alpha\beta\) additiv erzeugt; 
diese liegen im Bild der Polynomklasse von \(1,\alpha,\beta,\alpha\beta\).

\emph{(2) Relationen:} Aus \ref{squares-zero} folgt \(\alpha^2=\beta^2=0\) in \(H^\ast\); 
also ist \((\alpha^2,\beta^2)\subseteq\ker\Phi\).

\emph{(3) Keine weiteren Relationen:} Die Quotientenbasis \(\{1,\alpha,\beta,\alpha\beta\}\) 
bildet eine freie \(R\)-Basis der Quelle (als \(R\)-Modul). 
Da das Bild dieselbe Größe besitzt und \(\alpha\beta\) nach \ref{nonzero} nicht verschwindet, 
ist \(\ker\Phi= (\alpha^2,\beta^2)\).

\emph{(4) Ringstruktur:} Die graded-Kommutativität (Lemma \ref{graded-comm}) stimmt mit der Superregel 
in \(R[\alpha,\beta]/(\alpha^2,\beta^2)\) überein; damit ist \(\Phi\) ein Isomorphismus gradierter Ringe.
\end{PROOF}

\begin{REM}{Zusammenfassung als „freies Objekt modulo Dimension“}{summary}
Jede Sphäre hat \(H^\ast(S^n;R)\cong R[\gamma_n]/(\gamma_n^2)\) mit \(|\gamma_n|=n\).
Das Produkt \(S^p\times S^q\) liefert das \emph{Ring-Tensorprodukt} dieser beiden 
superkommutativen Algebren, also genau
\[
H^\ast(S^p;R)\,\widetilde{\otimes}\,H^\ast(S^q;R) 
\;\cong\; R[\alpha]/(\alpha^2)\;\widetilde{\otimes}\;R[\beta]/(\beta^2)
\;\cong\; R[\alpha,\beta]/(\alpha^2,\beta^2),
\]
wobei \(|\alpha|=p\), \(|\beta|=q\) und \(\alpha\beta=(-1)^{pq}\beta\alpha\).
\end{REM}


\pagebreak

\begin{DEF}[Koalgebra]
Dual zur Algebra ist eine \(R\)-\emph{Koalgebra} ein \(R\)-Modul \(C\)
mit einer Komultiplikation \(\Delta:C\to C\otimes_R C\)
und einer Kounität \(\varepsilon:C\to R\),
die koassoziativ und kounital sind:
\[
(\Delta\otimes\id)\circ\Delta=(\id\otimes\Delta)\circ\Delta,\qquad
(\varepsilon\otimes\id)\circ\Delta=\id=(\id\otimes\varepsilon)\circ\Delta.
\]
\end{DEF}

\begin{DEF}[Bialgebra und Hopf-Algebra]
Eine \emph{Bialgebra} über \(R\) ist zugleich Algebra und Koalgebra,
deren Strukturen kompatibel sind.
Besitzt sie zusätzlich einen Antipoden \(S:A\to A\)
(„inverse“ zur Komultiplikation),
so heißt sie \emph{Hopf-Algebra}.
\end{DEF}

\begin{REM}[Beispiele]
\begin{itemize}[leftmargin=*]
\item Die Gruppenalgebra \(R[G]\) ist eine Hopf-Algebra
mit \(\Delta(g)=g\otimes g\) und \(S(g)=g^{-1}\).
\item Der Kohomologiering \(H^\ast(G;R)\) einer topologischen Gruppe \(G\)
trägt eine natürliche \emph{Koalgebra}-Struktur über der Diagonalabbildung \(G\to G\times G\).
\end{itemize}
\end{REM}

\begin{CONC}[Gesamtstruktur]
Die bisher eingeführten Objekte lassen sich so anordnen:
\[
\text{Moduln} \;\xrightarrow{\text{Tensorprodukt}}\;
\text{Algebren}
\;\xrightarrow{\text{Dualisierung}}\;
\text{Koalgebren}
\;\xrightarrow{\text{Kombination}}\;
\text{Hopf-Algebren}.
\]
Algebren, Koalgebren und ihre graduierte oder super-Variante
bilden den strukturellen Rahmen nahezu aller algebraischen Theorien,
von Clifford- und Enveloping-Algebren bis zu Kohomologieringen.
\end{CONC}

% --- Kategorientheorie: Algebren als Monoidenobjekte in monoidalen Kategorien ---
\section{Kategorientheorie: Algebren als Monoidenobjekte in monoidalen Kategorien}
\label{sec:cat-monoids}

\begin{MOT}{Top-down Zielsetzung}{Cat-Algebra}
Wir formalisieren \emph{Algebra-Strukturen} als \emph{Monoidenobjekte} in einer passenden monoidalen Kategorie. Dadurch werden
(1) der tiefgreifende Zusammenhang zwischen Multiplikation und Tensorprodukt präzise,
(2) Unterschiede zwischen \(\Alg_R\) (alle \(R\)-Algebren) und \(\Alg_R^{\mathrm{gr}}\) (graduierte \(R\)-Algebren) konzeptionell klar.
\end{MOT}

\subsection{Monoidale Kategorien und Monoidenobjekte}

\begin{DEF}{Monoidale Kategorie}{monoidal}
Eine \emph{monoidale Kategorie} \((\mathcal{C},\otimes,\mathbb{1},\alpha,\lambda,\rho)\) besteht aus
\begin{itemize}[leftmargin=*]
  \item einem Bifunktor \(\otimes:\mathcal{C}\times\mathcal{C}\to\mathcal{C}\),
  \item einem Einheitsobjekt \(\mathbb{1}\in \mathcal{C}\),
  \item natürlichen Isomorphismen (Assoziator \(\alpha\) und Links-/Rechtseinheiten \(\lambda,\rho\)),
\end{itemize}
die die Mac~Lane-Kohärenzdiagramme (Pentagon, Dreieck) erfüllen.
\end{DEF}

\begin{DEF}{Monoidenobjekt}{monoid-object}
Ein \emph{Monoidenobjekt} in \(\mathcal{C}\) ist ein Objekt \(A\) mit Morphismen
\[
\mu:A\otimes A\to A,\qquad \eta:\mathbb{1}\to A,
\]
so dass 
\[
\mu\circ(\mu\otimes\id_A)\circ\alpha_{A,A,A} \;=\; \mu\circ(\id_A\otimes\mu),\qquad
\mu\circ(\eta\otimes\id_A)\circ\lambda_A^{-1}=\id_A=\mu\circ(\id_A\otimes\eta)\circ\rho_A^{-1}.
\]
\end{DEF}

\begin{REM}{Universelle Bilinearität}{bilin}
Ist \(\mathcal{C}=\Mod_R\) mit \(\otimes=\otimes_R\) und \(\mathbb{1}=R\), so ist jede \(R\)-bilineare Multiplikation
\(\mu:A\times A\to A\) äquivalent zu einem \emph{einzigen} \(R\)-linearen Morphismus
\(\bar\mu:A\otimes_R A\to A\) via der universellen Eigenschaft des Tensorprodukts. 
Damit sind \(R\)-Algebren genau Monoidenobjekte in \((\Mod_R,\otimes_R,R)\):
\[
\Alg_R \;\simeq\; \mathrm{Mon}\big(\Mod_R,\otimes_R,R\big).
\]
\end{REM}

\begin{EXA}{Adjunktion (frei \(\dashv\) Vergiss)}{adjunction}
Der Vergissfunktor \(U:\Alg_R\to\Mod_R\) besitzt einen linken Adjunkten \(T:\Mod_R\to\Alg_R\) (freie Algebra):
\[
\Hom_{\Alg_R}\!\big(T(V),A\big)\;\cong\;\Hom_{\Mod_R}\!\big(V,U(A)\big).
\]
\begin{center}
\begin{tikzcd}
\Mod_R \arrow[r, shift left=1.1ex, "T" ] & \Alg_R \arrow[l, shift left=1.1ex, "U" ] \\
\end{tikzcd}
\end{center}
Kategorisch: \(T\) ist der freie Monoidenfunktor in \(\Mod_R\).
\end{EXA}

\subsection{Der Zusammenhang \texorpdfstring{\(\mu\)}{mu} \emph{vs.} \texorpdfstring{\(\otimes\)}{tensor} (Punkt 1)}

\begin{CONC}{Multiplikation als Faltung über \(\otimes\)}{mu-vs-tensor}
\begin{enumerate}[leftmargin=*]
  \item \textbf{Linearisierung:} Bilinearität \(\Rightarrow\) eindeutiger \(\bar\mu:A\otimes A\to A\).
  \item \textbf{Kohärenz:} Assoziativität/Einheit von \(\mu\) sind genau die Monoidenaxiome relativ zu \((\otimes,\mathbb{1})\).
  \item \textbf{Functorialität:} Algebrahomomorphismen sind genau Monoidenmorphismen \((A,\mu,\eta)\to(B,\nu,\epsilon)\).
  \item \textbf{Symmetrien:} Ist \(\mathcal{C}\) \emph{symmetrisch} monoidal mit Symmetrie \(\sigma_{X,Y}:X\otimes Y\to Y\otimes X\),
        so ist \emph{(super-)Kommutativität} die Forderung \(\mu=\mu\circ\sigma_{A,A}\) (ggf. mit Vorzeichenregel).
\end{enumerate}
\end{CONC}

\begin{REM}{Tensorsicht auf Relationen}{quotients}
Freie Algebren \(T_R(V)=\bigoplus_{n\ge 0}V^{\otimes n}\) entstehen als freies Monoidenobjekt.
Symmetrische/äußere/Clifford-Algebren sind Quotienten von \(T_R(V)\) durch \(\otimes\)-kompatible Relationen; 
die Monoidenaxiome überleben den Quotienten genau dann, wenn das Relatorideal \(\otimes\)-stabil ist.
\end{REM}

\subsection{Symmetrische/Äußere/Clifford-Algebra als Monoiden-Quotienten}
\label{sec:monoid-quotients}




\begin{MOT}{Warum „als Monoiden-Quotient“?}{mono-quo}
Das Konzept \emph{„Algebra als Monoidenobjekt“} wird besonders klar, wenn man die
\emph{freie Algebra} \(T_R(V)\) als den \emph{freien Monoiden} in der Kategorie
\((\Mod_R,\otimes_R,R)\) betrachtet. Jede Relation, die wir in einer Algebra durchsetzen wollen,
ist dann nichts anderes als eine \emph{Monoidenrelation}, also eine Relation, die mit der
Multiplikationsstruktur verträglich ist.

\paragraph{Beispiel (Intuition mit \(R[x,y]\)):}
Betrachte den freien \(R\)-Modul \(V=R\{x,y\}\).
Die freie Algebra \(T_R(V)\) ist die direkte Summe aller Tensorpotenzen:
\[
T_R(V)=R \oplus V \oplus (V\otimes V)\oplus (V\otimes V\otimes V)\oplus\cdots
\]
und kann konkret als Menge aller endlichen „Wörter“ in den Symbolen \(x,y\) verstanden werden.
Die Multiplikation ist einfach die Konkatenation von Wörtern:
\[
(xy)\cdot(yx) = xy\,yx.
\]
Wenn wir aus dieser „nichtkommutativen“ Algebra \(R\langle x,y\rangle=T_R(V)\)
die kommutative Polynom\-algebra \(R[x,y]\) machen wollen,
müssen wir nur fordern, dass \(x\) und \(y\) vertauschbar sind:
\[
xy=yx.
\]
Diese Bedingung ist nichts anderes als eine \emph{Relation} zwischen den Elementen
\(x\otimes y\) und \(y\otimes x\) in \(V\otimes V\).  
Formal erzeugt man damit ein \emph{zweiseitiges Ideal}
\(\langle x\otimes y - y\otimes x\rangle\subseteq T_R(V)\), das von der Multiplikation
stabil bleibt (da auch jedes Produkt mit weiteren Symbolen die Relation bewahrt).  
Quotientiert man danach,
\[
R[x,y]\;=\;T_R(V)\big/\langle x\otimes y-y\otimes x\rangle,
\]
so ist dies genau der \emph{Monoiden-Quotient}, in dem die gewünschte Relation gilt.
\end{MOT}

\begin{DEF}{Monoid\-kongruenz und Idealkonstruktion}{monoid-ideal}
Formal sei \(j:V\to T_R(V)\) die kanonische Inklusion des \(R\)-Moduls \(V\)
in die Tensoralgebra und \(\iota_2=j\otimes j:V\otimes_R V\to T_R(V)\otimes_R T_R(V)\)
die Einbettung des Tensorquadrats.
Eine \(R\)-lineare Abbildung \(\rho:V\otimes_R V\to T_R(V)\)
definiert ein System von \emph{Multiplikationsrelationen}.  
Das von ihr erzeugte \emph{Monoidenideal} ist das kleinste zweiseitige Ideal
\(I\unlhd T_R(V)\), das alle Bilder \(\mathrm{im}(\rho)\) enthält und
unter der Multiplikation stabil ist:
\[
T_R(V)\cdot \mathrm{im}(\rho)\subseteq I,\qquad \mathrm{im}(\rho)\cdot T_R(V)\subseteq I.
\]
Der Quotient \(T_R(V)/I\) ist dann der \emph{Monoiden-Quotient} der freien Algebra,
in dem die durch \(\rho\) vorgegebenen Relationen gelten.
\end{DEF}

\begin{REM}{Koequalisierer-Sicht und universelle Beschreibung}{coeq}
In der monoidalen Kategorie \((\Mod_R,\otimes_R,R)\) kann man die Bildung eines
Monoiden-Quotienten auch rein kategorisch ausdrücken.
Man betrachtet dazu zwei parallele Morphismen, die „die alte“ und „die neue“ Multiplikation
realisieren. Für die Symmetrie \(\sigma:V\otimes V\to V\otimes V\) gilt:
\[
\begin{tikzcd}[column sep=huge]
V\otimes V \arrow[r, shift left=0.7ex, "\mu\circ\iota_2"]
               \arrow[r, shift right=0.7ex, swap, "\mu\circ\sigma\circ\iota_2"]
  & T_R(V) \arrow[r, two heads, "q"] & Q
\end{tikzcd}
\]
Das Objekt \(Q\) ist dann der \emph{Koequalisierer} dieses Morphismenpaares
in der Kategorie der \(R\)-Algebren.  
Konkret bedeutet dies: \(Q\) ist die „größte“ Quotienten\-Algebra von \(T_R(V)\),
in der die Gleichung \(vw=wv\) für alle \(v,w\in V\) gilt.
Somit kann man jede „relationale“ Algebra – etwa die symmetrische, äußere oder Clifford-Algebra –
als \emph{Koequalisiererdiagramm} in \(\Alg_R\) auffassen.
\end{REM}

\begin{EXA}{Das Beispiel \(R[x,y]\) nochmals kategorisch}{example-comm}
Das Diagramm
\[
\begin{tikzcd}[column sep=huge]
V\otimes V \arrow[r, shift left=0.7ex, "\mu\circ\iota_2"]
               \arrow[r, shift right=0.7ex, swap, "\mu\circ\sigma\circ\iota_2"]
  & T_R(V) \arrow[r, two heads] & R[x,y]
\end{tikzcd}
\]
besagt: \(R[x,y]\) ist der Koequalisierer, der die Relation \(xy=yx\) erzwingt.
\(\sigma\) vertauscht die Tensorfaktoren, und das Quotientieren bedeutet,
dass die Symmetrie der Tensoren zur Kommutativität der Multiplikation wird.
So wird die klassische Polynom\-algebra zum universellen Beispiel einer
\emph{kommutativen} Monoidenstruktur im Sinne der Kategorie \(\Mod_R\).
\end{EXA}

\paragraph{(1) Symmetrische Algebra als Monoiden-Quotient.}
\begin{DEF}{Symmetrische Algebra via Idealen}{sym-as-quot}
\[
S_R(V)\;=\;T_R(V)\big/\Big\langle\, j(v)j(w)-j(w)j(v)\;\bigm|\;v,w\in V\,\Big\rangle.
\]
Äquivalent: \(S_R(V)\) ist der Koequalisierer von \(\mu\circ\iota_2\) und \(\mu\circ\sigma\circ\iota_2\).
\end{DEF}

\begin{PROP}{Universelle Eigenschaft (frei-kommutativ)}{ue-sym}
Für jede \emph{kommutative} \(R\)-Algebra \(A\) und \(R\)-lineares \(f:V\to A\) existiert genau ein Algebrahomomorphismus \(\Phi:S_R(V)\to A\) mit \(\Phi\circ j=f\).
\end{PROP}

\paragraph{(2) Äußere Algebra als Monoiden-Quotient.}
\begin{DEF}{Äußere Algebra via Idealen}{ext-as-quot}
\[
\Lambda_R(V)\;=\;T_R(V)\big/\Big\langle\, j(v)j(v)\;\bigm|\;v\in V\,\Big\rangle.
\]
Über Ringen mit \(2\in R^\times\) ist dies äquivalent zum Quotienten nach 
\(\big\langle\, j(v)j(w)+j(w)j(v)\;\bigm|\;v,w\in V\,\big\rangle\),
da \(j(v+w)^2=0 \Rightarrow j(v)j(w)=-j(w)j(v)\).
\end{DEF}

\begin{PROP}{Universelle Eigenschaft (frei-alternierend)}{ue-ext}
Für jede \(R\)-Algebra \(A\) mit alternierender Multiplikation auf \(f(V)\) (\(f(v)^2=0\)) und \(R\)-lineares \(f:V\to A\) existiert genau ein Algebrahomomorphismus \(\Phi:\Lambda_R(V)\to A\) mit \(\Phi\circ j=f\).
\end{PROP}

\begin{REM}{Koequalisierer-Variante}{ext-coeq}
Ist \(2\in R^\times\), so erzwingt man Alternierung als Koequalisierer des Paares \(\mu\circ(\mathrm{id}+\sigma)\circ\iota_2\) und der Nullabbildung \(0\):
\[
\begin{tikzcd}[column sep=huge]
V\otimes V \arrow[r, shift left=0.7ex, "\mu\circ(\mathrm{id}+\sigma)\circ\iota_2"] \arrow[r, shift right=0.7ex, swap, "0"] & T_R(V) \arrow[r, two heads] & \Lambda_R(V).
\end{tikzcd}
\]
\end{REM}

\paragraph{(3) Clifford-Algebra als Monoiden-Quotient.}
\begin{DEF}{Clifford-Quotient}{cliff-as-quot}
Sei \(Q:V\to R\) eine quadratische Form (mit polarer Form \(B(v,w)=\tfrac{1}{2}(Q(v+w)-Q(v)-Q(w))\) falls \(2\in R^\times\)).
Dann
\[
\Cl(V,Q)\;=\;T_R(V)\big/\Big\langle\, j(v)j(v)-Q(v)\cdot 1 \;\bigm|\; v\in V\Big\rangle.
\]
Falls \(2\in R^\times\) ist dies äquivalent zu
\(
\big\langle\, j(v)j(w)+j(w)j(v)-2B(v,w)\cdot 1 \,\big\rangle.
\)
\end{DEF}

\begin{PROP}{Universelle Eigenschaft}{ue-cliff}
Für jede \(R\)-Algebra \(A\) und \(R\)-lineares \(f:V\to A\) mit \(f(v)^2=Q(v)\cdot 1_A\) existiert genau ein Algebrahomomorphismus \(\Phi:\Cl(V,Q)\to A\) mit \(\Phi\circ j=f\).
\end{PROP}




\begin{DEF}{Kan. Inklusion und Gradierung}{incl-grad}
Sei $R$ ein kommutativer Ring und $V$ ein $R$-Modul. Die Tensoralgebra
\[
T_R(V)=\bigoplus_{n\ge 0} V^{\otimes n},\qquad V^{\otimes 0}:=R
\]
ist natürlich $\mathbb{N}$-graduiert, mit Multiplikation durch Konkatenation:
$V^{\otimes i}\otimes V^{\otimes j}\to V^{\otimes(i+j)}$.
Die \emph{kanonische Inklusion}
\[
j:V\longrightarrow T_R(V)
\]
ist die Einbettung in Grad $1$ (identisch mit der direkten Summanden-Inklusion
$V\hookrightarrow R\oplus V\oplus V^{\otimes 2}\oplus\cdots$).
Für $v\in V$ schreiben wir oft schlicht $v$ statt $j(v)$.
\end{DEF}

\begin{DEF}{Quadratische Form und polare Form}{Q-B}
Eine \emph{quadratische Form} ist eine Abbildung $Q:V\to R$ mit
\begin{enumerate}[label=(\roman*),leftmargin=*]
\item $Q(rv)=r^2Q(v)$ für alle $r\in R$, $v\in V$,
\item die zugehörige \emph{polarisierte Form} $B:V\times V\to R$, definiert (falls $2\in R^\times$) durch
\[
B(v,w):=\tfrac12\big(Q(v+w)-Q(v)-Q(w)\big),
\]
ist $R$-bilinear und symmetrisch. Dann gilt $Q(v)=B(v,v)$.
\end{enumerate}
\textbf{Ohne} die Invertierbarkeit von $2$ arbeitet man primär mit einer \emph{symmetrischen bilinearen Form}
$B:V\times V\to R$ und setzt $Q(v):=B(v,v)$. In Aussagen, die explizit $Q$ verwenden, nehmen wir im Folgenden
$2\in R^\times$ an (Standardfall).
\end{DEF}

\begin{DEF}{„Grad-2“-Relatoren und das \emph{zweiseitige} Ideal}{twosided-ideal}
Sei $\rho:V\otimes_R V\to T_R(V)$ eine $R$-lineare Abbildung (Relator in Grad $2$).
Wir definieren die Menge der elementaren Relationen
\[
S_\rho\;:=\;\mathrm{im}(\rho)\;\subset\;T_R(V).
\]
Das \emph{von $S_\rho$ erzeugte zweiseitige Ideal} $I_\rho\unlhd T_R(V)$ ist
\[
I_\rho\;:=\;\sum_{a,b\in T_R(V)} a\cdot S_\rho\cdot b
\;=\;\Big\{\sum_k\,a_k\,s_k\,b_k\;\Bigm|\; s_k\in S_\rho,\;a_k,b_k\in T_R(V)\Big\}.
\]
\textbf{Wichtig:} Obwohl $S_\rho$ in Grad $2$ liegt, enthält $I_\rho$ Beiträge in \emph{allen} Graden:
Wenn $a\in V^{\otimes i}$, $s\in V^{\otimes 2}$, $b\in V^{\otimes j}$, dann
$a\,s\,b\in V^{\otimes(i+2+j)}$.
Ist $S_\rho$ homogen in Grad $2$, so ist $I_\rho$ ein \emph{homogenes (graduiertes) Ideal}.
\end{DEF}

\begin{REM}{Intuition an $V=R\{x,y\}$}{int-degree}
Setze $S_\mathrm{sym}=\{xy-yx\}$ (Kommutativität in Grad $2$).
Dann liegt z.\,B. $(xy-yx)z$ in Grad $3$ und gehört zu $I_\mathrm{sym}$;
ebenso $u(xy-yx)v$ für beliebige Wörter $u,v$.
Das Quotientieren $T(V)/I_\mathrm{sym}$ „propagiert“ also die Relation in alle Längen.
\end{REM}

\paragraph{Drei klassische Fälle (mit präzisen $\rho$, $I$).}

\begin{DEF}{Symmetrische Algebra als Monoiden-Quotient}{rho-sym}
\begin{align*}
\rho_{\mathrm{sym}} &: V\otimes_R V \longrightarrow T_R(V), &
\rho_{\mathrm{sym}}(v\otimes w) &:= j(v)j(w)-j(w)j(v).\\
I_{\mathrm{sym}} &:= \langle \mathrm{im}(\rho_{\mathrm{sym}})\rangle
= \Big\langle\, j(v)j(w)-j(w)j(v)\;\Bigm|\; v,w\in V\,\Big\rangle.
\end{align*}
Dann ist
\[
S_R(V):=T_R(V)\big/I_{\mathrm{sym}}
\]
die freie \emph{kommutative} $R$-Algebra auf $V$.
$I_{\mathrm{sym}}$ ist homogen (Grad $2$) und graduiert; die Gradierung von $T_R(V)$
descendiert auf $S_R(V)=\bigoplus_n S^n_R(V)$.
\end{DEF}

\begin{DEF}{Äußere Algebra als Monoiden-Quotient}{rho-ext}
\begin{align*}
\rho_{\mathrm{ext}} &: V\otimes_R V \longrightarrow T_R(V), &
\rho_{\mathrm{ext}}(v\otimes w) &:= j(v)j(w)+j(w)j(v).
\end{align*}
Das von $\mathrm{im}(\rho_{\mathrm{ext}})$ erzeugte zweiseitige Ideal
\[
I_{\mathrm{ext}}:=\Big\langle\, j(v)j(w)+j(w)j(v)\;\Bigm|\; v,w\in V\,\Big\rangle
\]
erzwingt Alternierung: Für $v=w$ ist $\rho_{\mathrm{ext}}(v\otimes v)=2\,j(v)^2$.
Ist $2\in R^\times$, so ist dies äquivalent zu $j(v)^2=0$.
Der Quotient
\[
\Lambda_R(V):=T_R(V)\big/I_{\mathrm{ext}}
\]
ist die äußere Algebra. Wieder ist $I_{\mathrm{ext}}$ homogen und graduiert.
\end{DEF}

\begin{DEF}{Clifford-Algebra als Monoiden-Quotient}{rho-cl}
Angenommen $2\in R^\times$ und $Q:V\to R$ ist quadratisch mit polarer Form $B$.
Definiere
\[
\rho_{\mathrm{Cl}}:V\otimes_R V\to T_R(V),\qquad
\rho_{\mathrm{Cl}}(v\otimes w)\;:=\;j(v)j(w)+j(w)j(v)\;-\;2\,B(v,w)\cdot 1.
\]
Dann
\[
I_{\mathrm{Cl}}:=\Big\langle\, j(v)j(w)+j(w)j(v)-2\,B(v,w)\cdot 1 \;\Bigm|\; v,w\in V\Big\rangle,
\qquad
\Cl(V,Q):=T_R(V)\big/I_{\mathrm{Cl}}.
\]
Für $v=w$ reduziert sich dies auf $2\,(j(v)^2-Q(v)\cdot 1)$, also $j(v)^2=Q(v)\cdot 1$.
Homogenität: $j(v)j(w)$ hat Grad $2$, die Konstante $1$ Grad $0$;
\emph{trotzdem} ist $I_{\mathrm{Cl}}$ ein graduiertes Ideal, denn die Relation ist
homogen, wenn man $1$ als Element in Grad $0$ zählt: Produkte $a\,[\cdots]\,b$
bleiben in wohldefinierten Graden $i+2+j$ bzw.\ $i+0+j$ für die skalaren Terme,
und die Gesamtheit der so erzeugten Elemente ist direkt eine gradweise Summe.
\end{DEF}

\begin{REM}{Alternative Formulierungen}{alt-forms}
\begin{itemize}[leftmargin=*]
\item Ohne Polarisationsannahme (kein $2^{-1}$) arbeitet man mit einer \emph{symmetrischen} Bilinearform
$B$ und fordert $j(v)^2=B(v,v)\,1$ sowie $j(v)j(w)+j(w)j(v)=(B(v,w)+B(w,v))\,1=2B(v,w)\,1$,
soweit sinnvoll. Viele Autoren definieren Clifford-Algebren primär aus $(V,B)$.
\item Für $\Lambda_R(V)$ kann man äquivalent $I=\langle j(v)\otimes j(v)\mid v\rangle$ (d.\,h. $v\wedge v=0$)
verwenden; zusammen mit der Assoziativität folgt dann $j(v)j(w)=-j(w)j(v)$.
\end{itemize}
\end{REM}

\begin{PROP}{Warum „Grad-2-Relationen“ alles bestimmen}{prop-spread}
Sei $S\subset V^{\otimes 2}\subset T_R(V)$ homogen in Grad $2$, $I=\langle S\rangle$ das zweiseitige Ideal.
Dann ist $I$ das direkte Summanden-weise (gradweise) von $T_R(V)$ erzeugte Ideal:
\[
I_n \;=\;\sum_{\substack{i+j+2=n\\ i,j\ge 0}} V^{\otimes i}\,S\,V^{\otimes j}
\quad\text{für alle }n\ge 2,\qquad I_0=I_1=0.
\]
Insbesondere „lebt“ $I$ nicht in Grad $2$ allein, sondern besitzt für alle $n\ge 2$ nichttriviale Komponenten.
\end{PROP}

\begin{EXA}{Konkretes $3$-Grad-Element im Kommutativitätsideal}{concrete3}
Für $V=R\{x,y,z\}$ und $S_\mathrm{sym}=\{xy-yx,\;xz-zx,\;yz-zy\}$ gilt z.\,B.
\[
(xy-yx)z \;\in\; I_{\mathrm{sym},3},\qquad
u(xz-zx)v \;\in\; I_{\mathrm{sym},\,\deg(u)+\deg(v)+2},
\]
für beliebige Wörter $u,v$. Der Quotient $T(V)/I_{\mathrm{sym}}$ identifiziert
alle Permutationen von Buchstaben in einem Wort.
\end{EXA}

\begin{REM}{Koequalisierer-Sicht (Grad-2 reicht als Generatoren)}{coeq-again}
Kategorisch erzwingen die parallelen Pfeile
\[
V\otimes V \;\rightrightarrows\; T_R(V),\qquad
v\otimes w \mapsto j(v)j(w),\;\; v\otimes w \mapsto j(w)j(v)
\]
(oder die Clifford-Version mit $2B(v,w)1$) via Koequalisierer genau dieselben Quotienten.
Die Tatsache, dass die Quelle $V\otimes V$ ist, spiegelt nur wider, dass wir \emph{Generatoren der Relationen}
in Grad $2$ platzieren; das resultierende Koequalisat ist dennoch eine globale Aussage über \emph{alle} Grade.
\end{REM}



\subsubsection*{Checkliste: „Ist mein Quotient ein Monoiden-Quotient?“}
\begin{itemize}[leftmargin=*]
\item \textbf{Homogenität:} (für Gradierungen) sind alle Relationen homogen? \(\Rightarrow\) die Gradstruktur descendiert.
\item \textbf{Zweiseitigkeit:} ist das von den Relationen erzeugte Ideal zweiseitig? \(\Rightarrow\) Multiplikation ist wohldefiniert.
\item \textbf{Stabilität unter \(\otimes\):} werden Relationen durch \(\mu\) und die monoidale Struktur respektiert? 
\item \textbf{Universalität:} ist die gewünschte universelle Eigenschaft (frei-kommutativ, frei-alternierend, Clifford) erfüllt?
\end{itemize}

\subsubsection*{Ein-Blick-Diagramme}
\begin{center}
\begin{tikzcd}[column sep=huge]
V\otimes V \arrow[r, shift left=0.7ex, "\mu\!\circ\!\iota_2"] \arrow[r, shift right=0.7ex, swap, "\mu\!\circ\!\sigma\!\circ\!\iota_2"] & T(V) \arrow[r, two heads, "q"] & S(V)
\end{tikzcd}

\medskip

\begin{tikzcd}[column sep=huge]
V\otimes V \arrow[r, shift left=0.7ex, "\mu\!\circ\!(\mathrm{id}+\sigma)\!\circ\!\iota_2"] \arrow[r, shift right=0.7ex, swap, "0"] & T(V) \arrow[r, two heads] & \Lambda(V) \quad (2\in R^\times)
\end{tikzcd}

\medskip

\begin{tikzcd}[column sep=huge]
V \arrow[r, "{v\mapsto \mu\!\circ\!\iota_2(v,v)}"] \arrow[dr, swap, "{v\mapsto Q(v)\cdot 1}"] & T(V) \arrow[d, two heads] \\
& \Cl(V,Q)
\end{tikzcd}
\end{center}

\subsection{Graduierte Algebren als Monoiden in \texorpdfstring{\(\Mod_R^{\mathbb{Z}}\)}{ModR^Z} (Punkt 2)}

\begin{DEF}{\(\mathbb{Z}\)-graduierte Kategorie}{graded-cat}
Sei \(\Mod_R^{\mathbb{Z}}\) die Kategorie der \(\mathbb{Z}\)-graduierten \(R\)-Moduln \(M=\bigoplus_{n\in\mathbb{Z}}M_n\)
mit grad-erhaltenden Morphismen. Monoidale Struktur:
\[
(M\otimes N)_k \;=\; \bigoplus_{i+j=k} M_i\otimes_R N_j,\qquad \mathbb{1}=R\text{ in Grad }0.
\]
Ist \(\Mod_R^{\mathbb{Z}}\) \emph{symmetrisch} monoidal (z.\,B. via Sign-Regel), spricht man von der \emph{super}-Kategorie.
\end{DEF}

\begin{DEF}{Graduierte Algebra als Monoidenobjekt}{graded-alg}
Eine \(\mathbb{Z}\)-\emph{graduierte \(R\)-Algebra} ist ein Monoidenobjekt \((A,\mu,\eta)\) in \((\Mod_R^{\mathbb{Z}},\otimes,\mathbb{1})\),
d.\,h.\ \(A=\bigoplus_{n\in\mathbb{Z}}A_n\) und \(\mu\) respektiert die Gradierung:
\[
\mu(A_i\otimes A_j)\subseteq A_{i+j}.
\]
\end{DEF}

\begin{PROP}{Vergleich \(\Alg_R\) vs. \(\Alg_R^{\mathrm{gr}}\)}{cmp-alg-graded}
\begin{enumerate}[leftmargin=*]
  \item \(\Alg_R \simeq \mathrm{Mon}(\Mod_R)\) erfasst \emph{alle} \(R\)-Algebren ohne zusätzliche Struktur.
  \item \(\Alg_R^{\mathrm{gr}} \simeq \mathrm{Mon}(\Mod_R^{\mathbb{Z}})\) erfasst Algebren \emph{mit} kompatibler Gradzerlegung.
  \item Der Vergissfunktor \(U:\Alg_R^{\mathrm{gr}}\to\Alg_R\) vergisst die Gradzerlegung; er ist \emph{konservativ}, aber nicht volltreu.
\end{enumerate}
\end{PROP}

\begin{REM}{Wann hat man \emph{nicht nur} eine Algebra, sondern eine graduierte?}{when-graded}
Äquivalente Situationen:
\begin{enumerate}[leftmargin=*]
  \item Es gibt eine direkte Summe \(A=\bigoplus_{n\in\mathbb{Z}}A_n\) mit \(\mu(A_i\otimes A_j)\subseteq A_{i+j}\).
  \item (Generator/Relationen) \(A\cong T_R(V)/\langle \mathcal{R}\rangle\) mit homogener Gradzuweisung \(\deg:V\to\mathbb{Z}\) und \emph{homogenen} Relationen \(\mathcal{R}\).
  \item (Gruppenaktion) Es gibt eine \emph{Gewichtsaktion} eines Monoid/Gruppe \(G\) (typisch \(G=\mathbb{Z}\) oder \(\mathbb{G}_m\) über Körper \(k\)) durch Algebra-Automorphismen, so dass
        \(A\) als direkte Summe der Gewichtsräume zerfällt; die Multiplikation addiert Gewichte.
  \item (Projektoren) Es existieren idempotente Projektoren \(\pi_n:A\to A\) mit \(A=\bigoplus_n \mathrm{im}(\pi_n)\) und \(\pi_{i+j}\circ\mu=\mu\circ(\pi_i\otimes\pi_j)\).
\end{enumerate}
Praktisch genügt: Weise Erzeugern Grade zu und verifiziere Homogenität aller Relationen.
\end{REM}

\begin{REM}{(Super-)Kommutativität in der graduierten Welt}{super-comm}
In einer \emph{symmetrisch} monoidalen \(\Mod_R^{\mathbb{Z}}\) mit \emph{Koszul-Zeichenregel} gilt für homogene \(a\in A_p,b\in A_q\):
\[
ab \;=\; (-1)^{pq}\,ba.
\]
Damit sind Kohomologieringe \(H^\ast(X;R)\) \emph{superkommutative} Monoidenobjekte.
\end{REM}

\subsection{Diagramme und Kohärenz}

\begin{EXA}{Assoziativitäts- und Einheitsdiagramme}{coherence}
\begin{center}
\begin{tikzcd}[column sep=large]
(A\otimes A)\otimes A \arrow[r, "\alpha_{A,A,A}"] \arrow[d, "\mu\otimes \id"'] 
& A\otimes(A\otimes A) \arrow[d, "\id\otimes \mu"] \\
A\otimes A \arrow[r, "\mu" ] & A
\end{tikzcd}

\bigskip

\begin{tikzcd}[column sep=huge]
\mathbb{1}\otimes A \arrow[r, "\lambda_A"] \arrow[d, "\eta\otimes \id_A"'] 
& A \\
A\otimes A \arrow[ur, "\mu"'] &
\end{tikzcd}
\end{center}
Die gleichen Diagramme gelten wortgleich in \(\Mod_R^{\mathbb{Z}}\), nur mit Gradverträglichkeit.
\end{EXA}

\subsection{Checkliste (praktisch)}

\begin{CONC}{Wie prüfe ich eine graduierte Algebra-Struktur?}{checklist}
\begin{enumerate}[leftmargin=*]
  \item \textbf{Gradwahl:} Lege \(\deg\) auf Erzeugern fest.
  \item \textbf{Homogene Relationen:} Überprüfe, dass jede Relation Summe gleichgradiger Terme ist.
  \item \textbf{Multiplikation:} Verifiziere \(\deg(xy)=\deg x+\deg y\) auf Erzeugern; schließe durch Bilinearität.
  \item \textbf{Einheit:} \(\deg 1_A=0\).
  \item \textbf{(Optional) Superkommutativ:} Prüfe die Vorzeichenregel im symmetrischen Setup.
\end{enumerate}
\end{CONC}

\subsection{Zusammenfassung der beiden Kernpunkte}

\begin{CONC}{Essenz}{essentials}
\begin{itemize}[leftmargin=*]
  \item \textbf{(1) Multiplikation \(\leftrightarrow\) Tensorprodukt:} Eine Algebra-Struktur ist \emph{genau} ein Monoidenobjekt relativ zu \((\otimes,\mathbb{1})\). Die Bilinearität von \(\mu\) wird vollständig durch die universelle Eigenschaft des Tensorprodukts modelliert.
  \item \textbf{(2) Ungraduiert vs. Graduiert:} \(\Alg_R\) betrachtet Monoiden in \(\Mod_R\); \(\Alg_R^{\mathrm{gr}}\) betrachtet Monoiden in \(\Mod_R^{\mathbb{Z}}\). Eine „bloße“ Algebra wird „graduiert“, sobald zusätzlich eine Gradzerlegung vorliegt, die mit \(\mu\) kompatibel ist (oder äquivalent: homogene Erzeuger/Relationen, Gewichtsaktion, Projektorzerlegung).
\end{itemize}
\end{CONC}



\pagebreak

\section{Operatoralgebren, Fockräume und (Anti-)Kommutationsrelationen}

\begin{DEF}{*-Algebra und Operatoralgebra}{op-alg}
Sei $\mathbb{K}\in\{\mathbb{R},\mathbb{C}\}$. 
Eine \emph{*-Algebra} $A$ über $\mathbb{K}$ ist eine assoziative Algebra mit einer
involutiven, antimultiplikativen Abbildung $*:A\to A$ (d.\,h.\ $(ab)^*=b^*a^*$, $(a^*)^*=a$).
Eine \emph{Operatoralgebra} (im engen Sinn) ist eine normierte *-Algebra von beschränkten Operatoren
$A\subset B(\mathcal{H})$ auf einem Hilbertraum $\mathcal{H}$, die unter $*$ abgeschlossen ist.
\begin{itemize}[leftmargin=*,nosep]
\item Eine \emph{$C^*$-Algebra} ist eine normierte *-Algebra mit $\|a^*a\|=\|a\|^2$, vollständig bzgl. der Norm.
\item Eine \emph{von-Neumann-Algebra} ist eine $*$-Unteralgebra von $B(\mathcal{H})$, die im starken (äquivalent: schwach-operator-) Abschluss abgeschlossen ist und die Identität enthält.
\end{itemize}
\end{DEF}

\begin{REM}{Physikalische Lesart}{phys-remark}
In der QFT/AQFT meint „Operatoralgebra“ typischerweise eine $C^*$- oder von-Neumann-Algebra von Observablen
(z.\,B.\ ein lokales Netz $\mathcal{O}\mapsto \mathcal{A}(\mathcal{O})$).
Konstruktiv startet man meist mit einer \emph{freien *-Algebra} auf Erzeugern (Felder, Moden, Einteilchenzustände)
und teilt dann durch das zweiseitige *-Ideal der \emph{Relationsgleichungen} (CCR/CAR, Feldgleichungen, Gauge-Constraints).
Die entstehende abstrakte *-Algebra besitzt dann je nach Kontext ein $C^*$- oder von-Neumann-Envelope.
\end{REM}

\subsection{Freie *-Algebra und Quotienten durch Relationen}

\begin{DEF}{Freie *-Algebra auf einem Vektorraum}{free-star}
Sei $E$ ein komplexer Vektorraum. Die \emph{freie *-Algebra} $F\langle E\rangle$ ist die freie assoziative Algebra
auf den Symbolen $E\cup E^\ast$ (formale Adjunkte), modulo den linearen Regeln
$(\alpha x+\beta y)^\ast=\bar\alpha x^\ast+\bar\beta y^\ast$ und $(xy)^\ast=y^\ast x^\ast$.
Jede lineare Abbildung $j:E\to A$ in eine *-Algebra $A$ mit gewähltem * hebt eindeutig zu einem *-Homomorphismus
$\Phi:F\langle E\rangle\to A$ an, der $j$ erweitert.
\end{DEF}

\begin{DEF}{Quotient durch Relationen}{relations}
Gegeben ein System von (polynomiellen) Relationen $\mathcal{R}\subset F\langle E\rangle$, sei
$I=\langle \mathcal{R}\rangle_{*}$ das kleinste zweiseitige *-Ideal, das $\mathcal{R}$ enthält.
Die \emph{abstrakte Operatoralgebra der Erzeuger mit Relationen $\mathcal{R}$} ist
\[
\mathcal{A}(E\mid \mathcal{R}) \;:=\; F\langle E\rangle\,\big/\,I.
\]
\end{DEF}

\subsection{Fockräume als symmetrische/äußere Algebren}

\begin{DEF}{Bosonischer/fermionischer Fockraum}{fock}
Sei $\mathfrak{h}$ ein komplexer Einteilchen-Hilbertraum.
\begin{align*}
\Gamma_s(\mathfrak{h}) &:= \bigoplus_{n=0}^\infty S^n \mathfrak{h} 
\quad\text{(bosonisch; symmetrische Tensormächte)},\\
\Gamma_a(\mathfrak{h}) &:= \bigoplus_{n=0}^\infty \Lambda^n \mathfrak{h}
\quad\text{(fermionisch; äußere Mächte)}.
\end{align*}
Topologisch sind dies die Hilbert-Vervollständigungen der symmetrischen bzw.\ äußeren Algebra über $\mathfrak{h}$.
\end{DEF}

\begin{DEF}{Erzeuger $a^*(f),a(f)$}{cre-ann}
Für $f\in \mathfrak{h}$ definiert man auf $\Gamma_s(\mathfrak{h})$ bzw.\ $\Gamma_a(\mathfrak{h})$
die \emph{Erzeugungs-} und \emph{Vernichtungsoperatoren} $a^*(f)$, $a(f)$ durch
\[
a^*(f):\; \psi^{(n)} \mapsto \mathrm{sym}(f\otimes \psi^{(n)}),\qquad
a(f):\; \psi^{(n)} \mapsto \text{(kontrahiere erste Komponente mit }f),
\]
wobei $\mathrm{sym}$ die (anti-)Symmetrisierung in der jeweiligen Kategorie ist.
\end{DEF}

\begin{PROP}{CCR und CAR}{ccr-car}
Mit dem Skalarprodukt $\langle\cdot,\cdot\rangle$ auf $\mathfrak{h}$ gelten:
\begin{enumerate}[label=(\roman*),leftmargin=*]
\item \textbf{Bosonen (CCR):}
\[
[a(f),a(g)] = 0,\quad [a^*(f),a^*(g)] = 0,\quad
[a(f),a^*(g)] = \langle f,g\rangle\,\mathbb{1}.
\]
\item \textbf{Fermionen (CAR):}
\[
\{a(f),a(g)\}=0,\quad \{a^*(f),a^*(g)\}=0,\quad
\{a(f),a^*(g)\} = \langle f,g\rangle\,\mathbb{1}.
\]
\end{enumerate}
\end{PROP}

\begin{REM}{Algebra aus Quotient}{quot-impl}
Die abstrakte CCR/CAR-\emph{*-Algebra} entsteht als Quotient der freien *-Algebra auf Symbolen
$\{a(f),a^*(f)\mid f\in\mathfrak{h}\}$ durch das *-Ideal, das von den obigen Relationen erzeugt wird.
Für CCR benutzt man oft die \emph{Weyl-Form} mit unitären $W(f)$ und
\[
W(f)W(g)=e^{-\,\tfrac{i}{2}\,\Im\langle f,g\rangle}\,W(f+g),\qquad W(f)^*=W(-f),
\]
und definiert die $C^*$-Algebra als die von $\{W(f)\}$ erzeugte $C^*$-Hülle.
\end{REM}

\subsection{Clifford-Algebra als CAR in Verkleidung}

\begin{DEF}{Clifford-Algebra}{cliff-def}
Sei $(V,B)$ ein reeller Prähilbertraum mit symmetrischer Bilinearform $B$.
Die \emph{Clifford-Algebra} $\Cl(V,B)$ ist der Quotient
\[
\Cl(V,B) \;=\; T(V)\,\big/\,\langle v\otimes v - B(v,v)\,1,\; v\in V\rangle,
\]
äquivalent durch $vw+wv=2B(v,w)\,1$.
\end{DEF}

\begin{PROP}{CAR $\Leftrightarrow$ Clifford}{car-cliff}
Betrachte $V$ als reellen Unterraum der komplexen Einteilchen-Hilberträume $\mathfrak{h}$ (vergiss die komplexe Struktur)
mit $B(v,w)=\Re\langle v,w\rangle$. Definiere auf der fermionischen Fockraum-Darstellung
\[
\psi(v)\;:=\; a^*(v)+a(v).
\]
Dann gilt
\[
\{\psi(v),\psi(w)\}\;=\;2\,B(v,w)\,\mathbb{1},
\]
und die durch $v\mapsto \psi(v)$ erzeugte *-Algebra ist isomorph zur komplexifizierten Clifford-Algebra
$\Cl(V,B)\otimes_\mathbb{R}\mathbb{C}$. In diesem präzisen Sinn sind die CAR die „operatorwertige“ Realisierung von
Clifford-Relationen.
\end{PROP}

\subsection{„Freie Theorie“ als freie/quotientenfreie Konstruktion}

\begin{CONC}{Freie QFT $\Rightarrow$ Fock aus Sym/Ext}{free-theory}
Die freie bosonische(fermionische) Theorie entsteht, indem man vom Einteilchenraum $\mathfrak{h}$
über den Funktor der zweiten Quantisierung
\[
\Gamma_s:\ \mathfrak{h}\mapsto \Gamma_s(\mathfrak{h}),\qquad
\Gamma_a:\ \mathfrak{h}\mapsto \Gamma_a(\mathfrak{h})
\]
den Fockraum (Vervollständigung der symmetrischen bzw.\ äußeren Algebra) bildet, darauf
Erzeuger/Vernichter definiert und die abstrakte Operatoralgebra als Quotient der freien *-Algebra
durch CCR bzw.\ CAR erhält. Feldoperatoren sind (nach Wahl geeigneter Testfunktionen/Moden) lineare
Kombinationen $a^*(f)\pm a(f)$ und werden als operatorwertige Distributionen modelliert.
\end{CONC}




\pagebreak

\section{Lie-Algebren und universelle Hülle}
\label{sec:lie-and-enveloping}

\begin{DEF}[Lie-Algebra]
Sei \(R\) ein kommutativer Ring mit \(1\). Eine \emph{Lie-Algebra} über \(R\) ist ein \(R\)-Modul \(\mathfrak{g}\) mit einer \(R\)-bilinearen Abbildung (Lie-Klammer)
\[
[-,-]:\mathfrak{g}\times\mathfrak{g}\to\mathfrak{g},
\]
so dass für alle \(x,y,z\in\mathfrak{g}\):
\begin{enumerate}[label=(\roman*), leftmargin=*]
\item \emph{Antisymmetrie:} \([x,y]= -[y,x]\).
\item \emph{Jacobi-Identität:} \([x,[y,z]]+[y,[z,x]]+[z,[x,y]]=0\).
\end{enumerate}
Ein Morphismus von Lie-Algebren ist ein \(R\)-linearer Homomorphismus \(\phi:\mathfrak{g}\to\mathfrak{h}\) mit \(\phi([x,y])=[\phi(x),\phi(y)]\).
\end{DEF}

\begin{EXA}[Kommutator aus Assoziativität]
Ist \(A\) eine assoziative \(R\)-Algebra, so wird \(A^{\operatorname{Lie}}\) durch den Kommutator
\([a,b]:=ab-ba\) zu einer Lie-Algebra. Dies definiert einen Funktor
\((-)^{\operatorname{Lie}}:\Alg_R\to\operatorname{Lie}_R\).
\end{EXA}

\begin{DEF}[Freie Lie-Algebra]
Für einen \(R\)-Modul \(V\) heißt \(\mathfrak{L}(V)\) eine \emph{freie Lie-Algebra} auf \(V\),
falls für jede Lie-Algebra \(\mathfrak{g}\) die universelle Eigenschaft gilt:
\[
\Hom_{\operatorname{Lie}_R}(\mathfrak{L}(V),\mathfrak{g}) \;\cong\; \Hom_{\Mod_R}(V,\mathfrak{g}).
\]
\emph{Konstruktion:} Schreibe \(T(V)=\bigoplus_{n\ge0}V^{\otimes n}\) für die Tensoralgebra. Dann ist
\[
\mathfrak{L}(V)\;=\; \text{die von }V\subset T(V)\text{ erzeugte Lie-Unteralgebra von }T(V)^{\operatorname{Lie}},
\]
wobei die Klammer in \(T(V)^{\operatorname{Lie}}\) der Kommutator \(ab-ba\) ist.
\end{DEF}

\begin{REM}[Freie Algebra impliziert Quotient auch im Lie-Fall]
Wie \(S(V),\ \Lambda(V)\) oder \(\Cl(V,Q)\) aus \(T(V)\) via Relationen entstehen, erhält man \(\mathfrak{L}(V)\) als \emph{freien} (nicht-assoziativen) Monoiden-Quotienten in der geeigneten (prä)Lie-Welt: Start ist stets \(T(V)\); Relationen (hier Jacobi/Antisymmetrie) werden als \(\otimes\)-verträgliche Kongruenzen implementiert.
\end{REM}

\begin{DEF}[Universelle assoziative Hülle]
Für eine Lie-Algebra \(\mathfrak{g}\) ist die \emph{universelle Hülle}
\[
U(\mathfrak{g})\;:=\; T(\mathfrak{g})\ \big/\
\left\langle x\otimes y - y\otimes x - [x,y]\ \middle|\ x,y\in\mathfrak{g}\right\rangle
\]
die assoziative \(R\)-Algebra, in der der Kommutator mit der Lie-Klammer übereinstimmt.
\end{DEF}

\begin{PROP}[Universelle Eigenschaft von \(U(\mathfrak{g})\)]
Für jede assoziative \(R\)-Algebra \(A\) und jeden Lie-Homomorphismus
\(\phi:\mathfrak{g}\to A^{\operatorname{Lie}}\) existiert genau ein eindeutiger
Algebrahomomorphismus \(\Phi:U(\mathfrak{g})\to A\) mit \(\Phi|_{\mathfrak{g}}=\phi\).
Damit ist \(U:\operatorname{Lie}_R\to\Alg_R\) linksadjungiert zu \((-)^{\operatorname{Lie}}:\Alg_R\to\operatorname{Lie}_R\).
\end{PROP}

\begin{PROP}[Poincaré–Birkhoff–Witt (PBW), skizziert]
Sei \(R\) ein Körper der Charakteristik \(0\) (oder hinreichend gut). Für eine Basis
\((x_i)_{i\in I}\) von \(\mathfrak{g}\) ist die Abbildung
\[
\operatorname{Sym}(\mathfrak{g}) \longrightarrow \operatorname{gr}\,U(\mathfrak{g})
\]
ein Isomorphismus von graduierten Algebren. Insbesondere bilden die „geordneten Monome“
\(x_{i_1}^{n_1}\cdots x_{i_k}^{n_k}\) (mit \(i_1<\dots<i_k\)) eine \(R\)-Basis von \(U(\mathfrak{g})\).
\end{PROP}

\begin{REM}[Bedeutung von PBW]
PBW zeigt, dass \(U(\mathfrak{g})\) „größenmäßig“ wie eine symmetrische Algebra aussieht
(\(\operatorname{gr}U\cong \operatorname{Sym}\)), während die Relationen \(xy-yx=[x,y]\) die
Lie-Struktur einkodieren. Das ist zentral für Darstellungen und Quantisierung.
\end{REM}

\begin{EXA}{\(\mathfrak{sl}_2\)}
Für \(\mathfrak{sl}_2(R)=\langle e,f,h\rangle\) mit
\([h,e]=2e,\ [h,f]=-2f,\ [e,f]=h\) ist
\[
U(\mathfrak{sl}_2)=T(e,f,h)\Big/\big\langle he-eh-2e,\ hf-fh+2f,\ ef-fe-h\big\rangle,
\]
und nach PBW bilden geordnete Monome \(f^a h^b e^c\) eine \(R\)-Basis.
\end{EXA}

\begin{EXA}[Heisenberg-Lie-Algebra]
Für Erzeuger \(p,q,z\) mit \([p,q]=z,\ [p,z]=[q,z]=0\) erhält man
\[
U(\mathfrak{h}) = T(p,q,z)\Big/\langle pq-qp-z,\ pz-zp,\ qz-zq\rangle,
\]
klassisches Modell für kanonische Vertauschungsrelationen (CCR) mit zentralem \(z\).
\end{EXA}

\begin{CONC}[Adjunktionsdreieck]
\[
\begin{tikzcd}[column sep=large]
\Mod_R \arrow[r, "T"] 
& \Alg_R \arrow[r, "(-)^{\operatorname{Lie}}"] \arrow[loop, distance=2em, in=55, out=125, "\text{assoziativ}"] 
& \operatorname{Lie}_R \arrow[l, bend left=20, "U"]
\end{tikzcd}
\]
\(\;\)Freie Objekte entstehen aus \(T\); der Übergang \(\Alg\to\operatorname{Lie}\) ist der Kommutator; \(U\) ist die Rückkehr in die assoziative Welt mit erzwungener Relation \(xy-yx=[x,y]\).
\end{CONC}

\begin{REM}[Physikalische Heuristik]
Poisson-Klammern \(\{\cdot,\cdot\}\) liefern Lie-Algebren von Observablen. Die Quantisierung ersetzt
\(\{A,B\}\) durch \((i\hbar)^{-1}[\,\hat A,\hat B\,]\) – formal ein Durchgang von einer
Lie-Struktur zur Operatoralgebra, modelliert durch \(U(\mathfrak{g})\) und Darstellungen.
\end{REM}




\pagebreak




\pagebreak


\printbibliography
\end{document}